%!TEX root = ../../main.tex
% ============================================================
% 物理/力学作图/八年级.tex
% 本文件由 main.tex 通过 \input 引入
% ============================================================


\section{解答：画出物体A所受的重力与物体A对墙壁的压力}

\vspace{0.6cm}

\begin{center}
\begin{tikzpicture}[>=Stealth, line join=round, line cap=round, scale=1.2]
  
  % ===================================
  % 一、 绘制原题干几何结构（黑色实线）
  % ===================================

  % 1. 绘制左侧竖直粗糙墙壁及其侧阴影修饰
  \foreach \y in {0.2, 0.5, ..., 4.0} {
    \draw[black, line width=0.6pt] (0, \y) -- (-0.3, \y-0.3);
  }
  \draw[black, line width=1.5pt] (0, 0) -- (0, 4.2);

  % 2. 绘制物体 A (置于墙面中央，方形)
  \draw[black, line width=1.5pt] (0, 1.1) rectangle (2.2, 3.1);
  \node at (1.1, 2.1) {\Large A};

  % 3. 绘制外部推力 F (水平向左，作用在物块右表面)
  \draw[->, black, line width=1.0pt] (4.0, 2.1) -- (2.2, 2.1) 
    node[midway, above, font=\large] {$F$};

  % ===================================
  % 二、 绘制题目所求的两个力（使用区分演示色）
  % ===================================

  % 【答案力 1】 物体A所受的重力 (Gravity)
  % 作用点：在物体 A 的几何中心，方向：竖直向下
  \coordinate (CenterA) at (1.1, 2.1);
  \fill[blue!80!black] (CenterA) circle (2pt);
  \draw[->, blue!80!black, line width=1.8pt] (CenterA) -- (1.1, 0.2) 
    node[right, font=\large\bfseries] {$G$};

  % 【答案力 2】 物体A对墙壁的压力 (Pressure on the wall)
  % 作用点：在物体与墙壁接触面的几何中心，即左侧垂直边中点
  % 方向：垂直于受力面（墙壁）向内（向左）
  \coordinate (WallContact) at (0, 2.1);
  \fill[red!80!black] (WallContact) circle (2pt);
  \draw[->, red!80!black, line width=1.8pt] (WallContact) -- (-1.8, 2.1) 
    node[above, font=\large\bfseries] {$F_{\text{压}}$};

\end{tikzpicture}
\end{center}

\vspace{0.6cm}

{\small\color{gray}
\textbf{解题说明：}\\
1. \textbf{确认重力要素：}地球表面所有物体都受到重力作用。由于物体 A 整体材质均匀且形状规则，其重心位于几何对角线的交点。重力的方向不受任何阻挡影响，永远处于“竖直向下”。\\
2. \textbf{确认压力要素：}“压力”是由相互挤压产生的相互作用力。题目要求画的是主体（物体A）对客体（墙面）施加的力。因此，压力的受力物体是“墙壁”，所以其“作用点”不能画在物体重心上，必须画在物体与墙面之间的实际接触面上（通常是接触截面的正中央）。因推力 $F$ 被挤压，压力方向垂直于受力墙面并向墙壁内部指向（水平向左）。
}

\vspace{0.6cm}

{\small\color{blue!70!black}
\textbf{绘图基本原理与步骤：}\\
\textbf{1. 物理环境拓扑：}利用 \texttt{rectangle} 以原点截距形式构筑标准的受力块；并通过 \texttt{foreach} 遍历向左下方平移 $-0.3$ 画出经典的左墙剖面线阴影。\texttt{line width=1.5pt} 加粗强化实体轮廓阻滞感。\\
\textbf{2. 力的三要素数学建模：}
\begin{itemize}
    \item \textbf{重力场：}锁定方块中心标 \texttt{coordinate}，绘制竖直向下垂线作为向量 $G$，并打上强调节点（\texttt{circle(2pt)}）。
    \item \textbf{受力面投影：}压力的实质物理根节点位于约束面上。必须切换原点至表面接触轴 \texttt{(0, y\_center)} 处作为力向量起点展开向左渲染；通过严格区分颜色的双线条区分物块内部力学属性以及墙体截面力学属性，大幅提升解析教学功能。
\end{itemize}
}


%% ==============================
%% 第6题：画出漂浮在液面上的木块所受力的示意图
%% ==============================


\section{解答：画出漂浮在液面上的木块所受力的示意图}

\vspace{0.6cm}

\begin{center}
\begin{tikzpicture}[>=Stealth, line join=round, line cap=round, scale=1.2]
  % 1. 绘制烧杯内的液体背景 (灰色填充)
  % 水面在 y=1.5，杯底在 y=0，边缘宽度为 1.8 到 -1.8，底部圆角 R=0.6 更自然
  \fill[gray!40] (-1.8, 1.5) -- (-1.8, 0.6) arc(180:270:0.6) -- (1.2, 0) arc(-90:0:0.6) -- (1.8, 1.5) -- cycle;
  
  % 2. 绘制液体表面线 (从杯壁到杯壁，中间部分会被木块遮挡)
  \draw[black, line width=0.8pt] (-1.8, 1.5) -- (1.8, 1.5);
  
  % 3. 绘制盛装液体的烧杯 (加粗黑线，自然平滑外翻，大圆弧过渡底角)
  \draw[black, line width=1.5pt] (-2.0, 2.5) to[out=-80, in=90] (-1.8, 1.5) -- (-1.8, 0.6) arc(180:270:0.6) -- (1.2, 0) arc(-90:0:0.6) -- (1.8, 1.5) to[out=90, in=260] (2.0, 2.5);
  
  % 4. 绘制漂浮的木块 (中间矩形，纯白填充遮挡后方液体线)
  % 木块中心设在 (0, 1.4)，木块大小为 1.2 * 1.0 (x = -0.6 to 0.6, y = 0.9 to 1.9)
  % 浸入液体深度为 1.5 - 0.9 = 0.6，露出液面为 1.9 - 1.5 = 0.4
  \draw[black, fill=white, line width=1.5pt] (-0.6, 0.9) rectangle (0.6, 1.9);

  % ===================================
  % 一、绘制题目所求的两个力（重力与浮力，等大反向）
  % ===================================
  
  % 标注重心
  \coordinate (Center) at (0, 1.4);
  \fill[black] (Center) circle (2pt);
  
  % 重力向下，加粗黑色矢量线 (长度 1.1)
  \draw[->, black, line width=1.8pt] (Center) -- (0, 0.3) 
    node[right, font=\large\bfseries] {$G$};
    
  % 浮力向上，与重力严格等长，加粗黑色矢量 (长度 1.1)
  \draw[->, black, line width=1.8pt] (Center) -- (0, 2.5) 
    node[right, font=\large\bfseries] {$F_{\text{浮}}$};

\end{tikzpicture}
\end{center}

\vspace{0.6cm}

{\small\color{gray}
\textbf{解题说明：}\\
1. \textbf{受力分析：}木块在水面处于“静止状态”（即漂浮），这是典型的二力平衡情形。因此木块整体上仅受到两个力的作用：一个是地球施加的竖直向下的重力 $G$，另一个是液体对其施加的竖直向上的浮力 $F_{\text{浮}}$。\\
2. \textbf{长度规格（画法关键）：}根据二力平衡约束条件，重力大小必须严格等于浮力大小（即 $|G| = |F_{\text{浮}}|$）。所以在实际作图阅卷中，这两条有向线段的端点长度\textbf{必须用直尺画得严格一样长}。\\
3. \textbf{作用点处理：}重力和浮力的作用点虽然物块受力实际微观部位不同（重力在重心，浮力在浸润体浮心），但为了作图教学方便、清晰表意宏观整体力学状态，通常将两个力共点且等效描点在物体的宏观几何重心位置上。并朝正上和正下方相互拉开对立力矩即可。
}

\vspace{0.6cm}

{\small\color{blue!70!black}
\textbf{绘图基本原理与步骤：}\\
\textbf{1. 容器渲染与层叠：}为还原最纯正的教辅原图样板，优先采取底部涂切逻辑 \texttt{\textbackslash fill[gray!40]} 渲染内壁包围的纯灰色区域作为液体区；直接横穿一根 $y=1.5$ 的水面实线；接着用 \texttt{\textbackslash draw[black]} 运用 \texttt{arc} 圆弧以及出入角向量 \texttt{out=-60, in=90} 等贝塞尔方法精密雕刻出极具自然张力的烧杯容器外侧翻边和杯筒。\\
\textbf{2. 绝妙的视差遮挡属性（Z-Index 运用）：}此处的木质矩形 \texttt{rectangle} 因为设定上就悬跨在浮潜截面线上，因此只需顺势为其添加上 \texttt{fill=white} 强力纯白填充遮罩属性，并将绘制层次（代码行数）在容器水体代码的下文处执行。系统渲染渲染器便会自动将该区块上方覆盖的基准水线、背景全副物理阻挡遮盖！轻松复活物理实体的三维感，省去庞大的手工切线代码。\\
\textbf{3. 二力等大的数学绑定：}重点在矢量表现端处，采用平行的固定距离数学位移坐标实现。源点坐标设在极简的原系对称中心 \texttt{(0, 1.4)} 上，那么拉升到 $y=2.5$ （长度 1.1）的力必与拉伸至下位 $y=0.3$ 的重力严格具有 1.1 的位移值。规避一切视觉透视歪斜或人眼误差，彻底合乎中学标量静定要求。
}

\vspace{1.0cm}

%% ==============================
%% 第7题：画出密度计在酒精中的大致位置
%% ==============================


\section{解答：画出密度计在酒精中的大致位置}

\vspace{0.6cm}

\begin{center}
\begin{tikzpicture}[>=Stealth, line join=round, line cap=round, scale=1.2]
  % ---------------------------------------------------
  % (a) 水中情形
  % ---------------------------------------------------
  \begin{scope}[shift={(-2.2, 0)}]
    % 1. 液体背景 (水，灰色填充)
    % 水面在 y=2.0，杯底在 y=0 
    \fill[gray!40] (-1.2, 2.0) -- (-1.2, 0.4) arc(180:270:0.4) -- (0.8, 0) arc(-90:0:0.4) -- (1.2, 2.0) -- cycle;
    
    % 2. 液体表面线 
    \draw[black, line width=0.8pt] (-1.2, 2.0) -- (1.2, 2.0);
    
    % 3. 烧杯轮廓
    \draw[black, line width=1.5pt] (-1.4, 2.5) to[out=-80, in=90] (-1.2, 1.8) -- (-1.2, 0.4) arc(180:270:0.4) -- (0.8, 0) arc(-90:0:0.4) -- (1.2, 1.8) to[out=90, in=260] (1.4, 2.5);
    
    % 4. 密度计 (在水中)
    % 假设水深 2.0，设浸入深度 h1 = 1.2，则底端在 y=0.8
    % 本体使用纯白填充以遮挡背后的水位实线
    \draw[black, fill=white, line width=1.2pt] (-0.2, 0.8) rectangle (0.2, 4.3);
    
    % 密度计内部配重物（底部铅粒/铁丝示意）
    \fill[gray!60] (-0.15, 0.85) rectangle (0.15, 1.1);
    \draw[black, line width=0.6pt] (-0.08, 1.1) to[out=60, in=-120] (0.08, 1.4) to[out=60, in=-120] (-0.08, 1.7) to[out=60, in=-120] (0.0, 2.0);
    
    % 5. 标注
    \draw[black, line width=0.8pt] (1.3, 0.8) -- (2.0, 0.8) node[right, font=\large] {水};
    \node at (0, -0.8) {\Large (a)};
  \end{scope}

  % ---------------------------------------------------
  % (b) 酒精中情形 (作图答案)
  % ---------------------------------------------------
  \begin{scope}[shift={(2.2, 0)}]
    % 1. 液体背景 (酒精，灰色填充)
    \fill[gray!40] (-1.2, 2.0) -- (-1.2, 0.4) arc(180:270:0.4) -- (0.8, 0) arc(-90:0:0.4) -- (1.2, 2.0) -- cycle;
    
    % 2. 液体表面线 
    \draw[black, line width=0.8pt] (-1.2, 2.0) -- (1.2, 2.0);
    
    % 3. 烧杯轮廓
    \draw[black, line width=1.5pt] (-1.4, 2.5) to[out=-80, in=90] (-1.2, 1.8) -- (-1.2, 0.4) arc(180:270:0.4) -- (0.8, 0) arc(-90:0:0.4) -- (1.2, 1.8) to[out=90, in=260] (1.4, 2.5);
    
    % 4. 密度计 (在酒精中，使用稍粗黑线作为本题学生需绘制的作图答案)
    % 设酒精的密度约为水的 0.8 倍，由于整体要求浮力一致（均等于重力），排开液体体积必须成反比增大。
    % 所以在酒精中浸入深度 h2 = h1 * (1.0/0.8) = 1.25 * 1.2 = 1.5。即浸入深度加深 0.3。
    % 则密度计底端坐标应该设定为 y = 2.0 - 1.5 = 0.5。对应顶端下降为 y = 4.0。
    \draw[black, fill=white, line width=1.8pt] (-0.2, 0.5) rectangle (0.2, 4.0);
    
    % 密度计内部配重物同步下移
    \fill[gray!60] (-0.15, 0.55) rectangle (0.15, 0.8);
    \draw[black, line width=0.6pt] (-0.08, 0.8) to[out=60, in=-120] (0.08, 1.1) to[out=60, in=-120] (-0.08, 1.4) to[out=60, in=-120] (0.0, 1.7);
    
    % 5. 标注
    \draw[black, line width=0.8pt] (1.3, 0.8) -- (2.0, 0.8) node[right, font=\large] {酒精};
    \node at (0, -0.8) {\Large (b)};
  \end{scope}

\end{tikzpicture}
\end{center}

\vspace{0.6cm}

{\small\color{gray}
\textbf{解题说明：}\\
1. \textbf{浮沉条件与浮力常性：}简易密度计无论是在水中还是在酒精中，都能够处于稳定的“漂浮状态”。根据力学平衡的静力学方程可以得知：漂浮情况下其受到的竖直向上的浮力 $F_{\text{浮}}$ 必须等于其自身的重力 $G$。即 $F_{\text{浮水}} = F_{\text{浮酒}} = G$。\\
2. \textbf{排开体积与液体密度的反比例解析：}根据阿基米德浮力公式 $F_{\text{浮}} = \rho_{\text{液}} g V_{\text{排}} = \rho_{\text{液}} g S h_{\text{浸}}$。因为纯酒精的密度小于普通水的密度（$\rho_{\text{酒精}} < \rho_{\text{水}}$），在保证产生的浮力大小相同、自身横截面积 $S$ 全局固定的前提下，密度计在酒精中排开这种液体的体积刻度必须强行增大补偿，即在酒精中浸没的深度 $h_{\text{浸}}$ 必定比在水中更深。\\
3. \textbf{作图结论：}因此在绘制图 (b) 酒精中的密度计时，不仅形状要保持完全一致，最关键的考察点在于它的上下坐标应当集体“下沉”。其浸入酒精液面下的细长部分，要远深于浸入水中的长度。
}

\vspace{0.6cm}

{\small\color{blue!70!black}
\textbf{绘图基本原理与步骤：}\\
\textbf{1. 环境等比例封装与镜像复用：}由于 (a) 图和 (b) 图属于同一规格下截然不同的物理容器态，利用 TikZ 内置的局部 \texttt{scope} 嵌套配合 \texttt{shift=\{(-2.2, 0)\}} 以及 \texttt{shift=\{(2.2, 0)\}} 横移逻辑，将左右图分块平移组合。彻底避免了多容器情况下的复杂坐标重构。\\
\textbf{2. 密度计管身的遮挡图层处理：}如上题般如出一辙地，密度计的身管主体 \texttt{rectangle} 参数中加入了 \texttt{fill=white}。这简单粗暴却极具效果地将后方横移截断穿梭的水面波浪实线遮掩，严格契合水面与漂浮物体之间的真实“三维相干与物理透视阻挡”。另外由于密度计量器底部通常缠绕有细铁丝或者配重小铅球以保证直立，我运用 \texttt{rectangle} 和多重交替的贝塞尔方程模拟了其细铁丝纹路（水和酒精中作了随动的相对下移复原），大大增强了还原考试配图的教辅正规感。\\
\textbf{3. 高精度的浮力下沉参数化：}在定量上并不通过任意瞎比划来体现“下沉”：先验地设定了水深 $H=2.0$ 相对参照，刻画其在水内浸没 $h_1 = 1.2$。严格依据高中物理定则取 $\rho_{\text{酒}} \approx 0.8\text{ g/cm}^3$，由此确凿算出如果等浮力该计身在酒精中必会下沉至深度 $h_2 = h_1 \times (1.0 \div 0.8) = 1.5$ 处。从而代码直接敲定了在酒精中该柱体底部的局部标高必然精准停留在 $y = 2.0 - 1.5 = 0.5$；这套严谨的数字逻辑完全避免了凭空人眼主观定位失察的问题，真正用代码逻辑“算出了结果图”而非仅为图做图！
}

\vspace{1.0cm}

%% ==============================
%% 第8题：画出浸没在水中的小球所受浮力的示意图
%% ==============================


\section{解答：画出浸没在水中的小球所受浮力的示意图}

\vspace{0.6cm}

\begin{center}
\begin{tikzpicture}[>=Stealth, line join=round, line cap=round, scale=1.2]
  % 1. 绘制水面 (实线)
  \draw[black, line width=1.0pt] (-2.5, 2.5) -- (2.5, 2.5);

  % 2. 绘制水的虚线背景 (表示水)
  \begin{scope}
    % 利用 clip 限制虚线不出界
    \clip (-2.5, 0) rectangle (2.5, 2.5);
    % 奇数行虚线
    \foreach \y in {0.2, 0.6, ..., 2.4} {
      \draw[black, dashed, dash pattern=on 4pt off 4pt, line width=0.4pt] (-2.5, \y) -- (2.5, \y);
    }
    % 偶数行错开相位虚线，这更符合典型的流体题干画风
    \foreach \y in {0.4, 0.8, ..., 2.4} {
      \draw[black, dashed, dash pattern=on 4pt off 4pt, dash phase=4pt, line width=0.4pt] (-2.5, \y) -- (2.5, \y);
    }
  \end{scope}

  % 3. 绘制完全浸没的小球 (实心白底黑边圆形)
  % 为了不让水的虚线穿透小球，这里为小球填充了不透明纯白色（fill=white）
  \coordinate (BallCenter) at (0, 1.2);
  \draw[black, fill=white, line width=1.2pt] (BallCenter) circle (0.5);

  % ===================================
  % 一、 绘制题目所求的答案力（浮力）
  % ===================================
  
  % 标注重心
  \fill[black] (BallCenter) circle (2pt);
  
  % 浸没在水中的物体受到的浮力方向竖直向上，作用点在重力几何中心
  \draw[->, black, line width=1.8pt] (BallCenter) -- (0, 3.2) 
    node[right, font=\large\bfseries] {$F_{\text{浮}}$};

\end{tikzpicture}
\end{center}

\vspace{0.6cm}

{\small\color{gray}
\textbf{解题说明：}\\
1. \textbf{确认受力物体与浮力三要素：}题目要求画出“小球”所受“浮力”的示意图。浮力产生的原因是液体对物体上下表面的压力差，根据阿基米德原理及其受力特性，物理学中规定浮力的\textbf{方向始终是竖直向上}的。\\
2. \textbf{确定作用点：}浮力的等效作用点称为“浮心”，对于材质均匀且完全浸没的规则形体（如本题中的小球），浮心与其几何中心（重心）重合。因此，我们将画笔的起点（受力点）点在小球的圆心处。\\
3. \textbf{作图规范：}从圆心开始，沿着竖直向上的方向画一条带实心箭头的线段，表示浮力向量，并在箭头端部附近标注大写字母符号 $F_{\text{浮}}$ 即可完成规范解答。
}

\vspace{0.6cm}

{\small\color{blue!70!black}
\textbf{绘图基本原理与步骤：}\\
\textbf{1. 水域环境流体构建：}基于顶部的 \texttt{draw} 绘制水面实线界限；随之通过 TikZ 的循环语法 \texttt{\textbackslash foreach} 分奇数和偶数阵列配合 \texttt{dashed} 和错位配置 \texttt{dash phase=4pt} 参数，快速铺展出等间距且交错的精美点阵横向虚线群，高度还原了物理原题教辅中典型的纯正标量“水池”虚线质感。\\
\textbf{2. 实体遮挡技巧：}利用 \texttt{\textbackslash draw[fill=white]} 语法绘制小球。这不但能够勾勒出球体主边界 \texttt{circle(0.5)}，更关键的是白底不透明度能在渲染层级上自然“强势掩盖遮挡”下方的水纹穿梭虚线群，从而彻底避免了水层虚假穿透球体内部的低级失真视觉错位状况。\\
\textbf{3. 力学标定：}单独抽出统一的作用点几何坐标 \texttt{coordinate (BallCenter)} 声明。直接在球心用 \texttt{fill} 打上实心黑点标记重心，并在视觉的顶层引出加宽加粗并附带 \texttt{Stealth} 实心箭头的竖直向上长力矩射线，契合标准作图题答案的高精读样张设计。
}

\vspace{1.0cm}

%% ==============================
%% 第9题：画出试管转到竖直位置时的液面
%% ==============================


\section{解答：画出试管转到竖直位置时的液面}

\vspace{0.6cm}

\begin{center}
\begin{tikzpicture}[>=Stealth, line join=round, line cap=round, scale=1.2]
  
  % ===================================
  % 一、 绘制原题干几何结构（黑色实线与虚线）
  % ===================================

  % 1. 绘制水平地面及阴影
  \draw[black, line width=1.0pt] (-2.0, 0) -- (4.0, 0);
  \foreach \x in {-1.8, -1.4, ..., 3.8} {
    \draw[black, line width=0.6pt] (\x, 0) -- (\x-0.2, -0.2);
  }

  % 2. 绘制位置 1 的试管（倾斜）和内部初始水
  \begin{scope}[shift={(0.35, 0.25)}, rotate=-40]
    % 填充位置 1 的水 (保持水平)
    % 水柱中心在 L=2.0 处。全局水平面在局部表现为斜率为 tan(40)
    \fill[gray!50] (-0.25, {2.0 - 0.25*tan(40)}) -- (-0.25, 0) arc(-180:0:0.25) -- (0.25, {2.0 + 0.25*tan(40)}) -- cycle;
    % 水面轮廓线
    \draw[black, line width=1.0pt] (-0.25, {2.0 - 0.25*tan(40)}) -- (0.25, {2.0 + 0.25*tan(40)}); 
    
    % 绘制位置 1 的试管外壳 (实线)
    \draw[black, line width=1.2pt] (-0.25, 4.0) -- (-0.25, 0) arc(-180:0:0.25) -- (0.25, 4.0) -- cycle;
  \end{scope}

  % 标注"1"：处于试管1右侧水面附近
  \node at (2.6, 2) {\Large 1};

  % 3. 绘制位置 2 的试管（竖直，虚线）
  \begin{scope}[shift={(0, 0.25)}]
    % ===================================
    % 二、 绘制题目所求的答案液面部分
    % ===================================
    % 倾斜时水的中心线在 L=2.0，竖直时其高度仍为 L=2.0。此时液面高度必定高于倾斜时。
    \fill[gray!50] (-0.25, 2.0) -- (-0.25, 0) arc(-180:0:0.25) -- (0.25, 2.0) -- cycle;
    
    % 加粗强调竖直时的液面（本题作图答案），并保持原题风格，使用明显稍粗的黑线
    \draw[black, line width=2.0pt] (-0.25, 2.0) -- (0.25, 2.0);
    
    % 绘制试管轮廓 (虚线)
    \draw[dashed, black, line width=1.0pt] (-0.25, 4.0) -- (-0.25, 0) arc(-180:0:0.25) -- (0.25, 4.0) -- cycle;
    
    \node at (0.4, 4.4) {\Large 2};
  \end{scope}

  % 4. 添加转动标志（弯曲箭头）
  \draw[->, black, line width=1.0pt] (2.0, 3.6) to[out=130, in=20] (0.5, 4.2);

\end{tikzpicture}
\end{center}

\vspace{0.6cm}

{\small\color{gray}
\textbf{解题说明：}\\
1. \textbf{作图分析（画出液面位置）：}试管内水的质量和体积是不变的，设试管的底面积为 $S$。当试管倾斜搁置在位置 1 时，水面是水平的，由于圆柱斜截体体积定理，它等效于将水柱“拉平”后齐平在其几何中心处。设这管水若竖直拉平时的管内长度为 $L$。因此，当试管立起到位置 2 时，液柱沿管壁长度恰好是原来位置 1 液面中心的所在管壁长度（即中心点仍然在内部坐标 $L$ 处）。由于试管位置 2 完全竖直，虽然水在管内所占的“长度刻度”不变，但其中心相对于桌面的\textbf{竖直高度相对明显增高}（原高度等于所在管壁长度乘以倾角对应的正弦或余弦分量），直观表现为：图上该水面直线相较于位置 1 水面的中央具有明显的直观上升。\\
2. \textbf{填空（压强变化）：}填空答案为\textbf{变大}。根据液体压强公式 $p = \rho g h$，其中 $h$ 为液体深度（液面到管底的垂直高度）。在此管由斜起立倒直转动的过程中，其内水体也整体被抬起，最高面（液面）对最低面（试管底）的垂直直线深度 $h$ 明显增大。因此，水对试管底部的压强\textbf{变大}。
}

\vspace{0.6cm}

{\small\color{blue!70!black}
\textbf{绘图基本原理与步骤：}\\
\textbf{1. 结构复刻：}利用 TikZ 的 \texttt{foreach} 遍历坐标绘制等距斜短线阵列，构建带阴影的基准水平面环境。\\
\textbf{2. 旋转坐标系与切面一致性质保障：}为确保倾斜试管无论旋转多大角度，底部的半圆极点始终“悬挂并相切”于桌面，不能简单依赖从管口等处向下方旋转探测位置。需要将局部 \texttt{scope} 的中心原点（旋转锚点）精准设置在距离地面高度正好等于试管圆倒角半径 $R=0.25$ 的位置：\texttt{shift=\{(x\_c, 0.25)\}}。以此点为锚固圆心执行 \texttt{rotate=-40} 的切向倾角自变换后，绘制出底面半圆曲线的物理最低点依然始终严格切附在地面桌上，此思路规避了手工直接的推算复杂切面坐标的盲点。\\
\textbf{3. 水位倾斜的自适应建模与守恒体积体现：}在自身处于倾斜姿态的局部坐标系 \texttt{scope} 内部视图下，原本全局的重力“水平面”变成了有确定斜率的倾斜直线，斜率为全局转角即 pgfmath 解析值 \texttt{tan(40)}。基于流体特性，我们只需固定体积约束下常量管中水长 $L=2.0$ 作为动态水面的纵截面中心点集；直接借助公式表达式 \texttt{y = 2.0 $\pm$ R*tan(40)} 便能精确求取倾斜水面的左右两侧切点相对坐标并执行安全闭合 \texttt{cycle} 进行液域着色；对应于受力终态也就是位置 2 （完全竖直），出于体积等价等效公理，该终态液面高度可以直接直接的切落截断于相对常量 \texttt{y=2.0} 。辅以加重加粗的纯黑黑实线，这种直接渲染能够使该体积逻辑高度无论在参数变化中均稳定精确大于倾斜状态时的物理实际视觉中心标高，真正实现在原汁原味代码化和物理规律的双重自洽检验。
}

%% ==============================
%% 第18题：画出下列各力的示意图
%% ==============================


\section{解答：画出下列各力的示意图}

\vspace{0.6cm}

\begin{center}
\begin{tikzpicture}[>=Stealth, line join=round, line cap=round, scale=1.0]

  % ---------------------------------------------------
  % (1) 图(a)：用 40N 力拉小车
  % ---------------------------------------------------
  \begin{scope}[shift={(-3.5, 0)}]
    % 1. 绘制水平地面及阴影
    \draw[black, line width=1.0pt] (-2.5, 0) -- (2.5, 0);
    \foreach \x in {-2.3, -2.1, ..., 2.3} {
      \draw[black, line width=0.6pt] (\x, 0) -- (\x-0.15, -0.15);
    }
    
    % 2. 绘制小车车体 (矩形主体)
    \draw[black, line width=1.2pt, fill=white] (-1.0, 0.4) rectangle (1.0, 1.8);
    % 小车右侧的连接挂点 (半圆环)
    %\draw[black, line width=1.2pt] (1.0, 1.3) arc(90:-90:0.2);
    
    % 3. 绘制车轮
    \draw[black, line width=1.2pt, fill=white] (-0.5, 0.2) circle (0.2);
    \draw[black, line width=1.2pt, fill=white] (0.5, 0.2) circle (0.2);
    
    % 4. 绘制答案：拉力 F 的示意图
    % 作用点在右侧挂钩中心
    \coordinate (HookA) at (1, 1.1);
    \fill[black] (HookA) circle (2pt);
    
    % 拉力向量，大小40N，方向与水平成30度右上
    \draw[->, black, line width=1.8pt] (HookA) -- ++(30:2.5) 
      node[right, font=\large\bfseries] {$F = 40\,\text{N}$};
    
    % 辅助虚线与角度标注
    \draw[dashed, black, line width=0.8pt] (HookA) -- ++(1.8, 0);
    \draw[black, line width=0.8pt] (HookA) ++(1.0, 0) arc(0:30:1.0);
    \node[font=\small] at ($(HookA) + (15:1.35)$) {$30^\circ$};
    
    % 图示序号
    \node at (0, -0.8) {\Large (a)};
  \end{scope}

  % ---------------------------------------------------
  % (2) 图(b)：细线下小球受到的重力
  % ---------------------------------------------------
  \begin{scope}[shift={(5.5, 0)}]
    % 1. 绘制水平底座
    \draw[black, fill=gray!80, line width=1pt] (-2.5, 0) rectangle (2.5, 0.3);
    
    % 2. 绘制右侧的垫块支撑
    \draw[black, fill=gray!30, line width=1pt] (0.8, 0.3) rectangle (1.7, {0.3 + 3.0*tan(15)});
    
    % 3. 绘制倾斜的斜面木板与 L 型支架
    \begin{scope}[shift={(-2.2, 0.4)}, rotate=15]
      % 斜面木板
      \draw[black, fill=white!90!gray, line width=1pt] (0, -0.1) rectangle (4.2, 0.2);
      % L型支架 (垂直固定在斜面上)
      \draw[black, line width=2pt] (3.2, 0.2) -- (3.2, 3.2) -- (1.5, 3.2);
      % 支架悬挂点
      \coordinate (HookB) at (1.5, 3.2);
      
      % 支架顶端垂下的那条垂直于斜面的参考虚线 (法线)
      \draw[dashed, black, line width=0.8pt] (HookB) -- (1.5, 0.5);
    \end{scope}
    
    % 4. 绘制原题干悬挂的细线与小球 (悬垂在全局重力方向，真实垂直)
    % 从挂点全局向下绘制细线
    \coordinate (BallB) at ($(HookB) + (0, -1.8)$);
    \draw[black, line width=1.0pt] (HookB) -- (BallB);
    
    % 绘制带灰度渐变立体感的小球
    \shade[ball color=gray!40] (BallB) circle (0.3);
    \draw[black, line width=0.6pt] (BallB) circle (0.3);
    
    % 5. 绘制答案：小球受到的重力 G 的示意图
    \fill[black] (BallB) circle (2pt);
    % 重力始终竖直向下。质量 0.5kg，即 G = mg = 0.5kg * 9.8N/kg = 4.9N (近似 5N)
    \draw[->, black, line width=1.8pt] (BallB) -- ++(0, -1.6) 
      node[right, font=\large\bfseries] {$G = 5\text{N}$};
    
    % 图示序号
    \node at (0, -0.8) {\Large (b)};
  \end{scope}

\end{tikzpicture}
\end{center}

\vspace{0.6cm}

{\small\color{gray}
\textbf{解题说明：}\\
1. \textbf{图(a) 拉力作图分析：}拉力 $F$ 是小车受到的约束物理力，作用点应该画在小车右侧的挂钩节点上。题干明确指出拉力“与水平方向成 $30^\circ$ 角”、“斜向右上方”且“大小为 $40\text{ N}$”。作图时严循参数，从挂钩中心出发作水平线作基准，用虚线引出。用弧线偏转 $30^\circ$ 以确定方向画出一条带箭头的实心力长线段，末端精准定标并标注表达式 $F = 40\,\text{N}$。\\
2. \textbf{图(b) 重力作图分析：}重力 $G$ 来源于地球的万有引力吸引。无论木板支持面是如何倾斜（如支架），**重力的方向永远是竖直向下的**。恰巧静止的小球那根吊线正好天然揭示了它的指向。题图中那根倾斜的虚线是支架垂线（即斜面的法线代表物），千万不可画错其上。准确的作图必须从小球中心圆心出发，向下严格垂直向地板画出并带出剪头。结合题目初算重力大小为 $G = mg = 0.5\,\text{kg} \times 9.8\,\text{N/kg} = 4.9\,\text{N}$（取 $g=10\,\text{N/kg}$ 则为 $5\text{ N}$），将这一数值准确附在箭头边即完成了标准的示意图作图。
}

\vspace{0.6cm}

{\small\color{blue!70!black}
\textbf{绘图基本原理与步骤：}\\
\textbf{1. 极坐标及数学库的应用：}在题(a)对于拉力 $F$ 角的作图中，极客般地利用了 TikZ 的局部相对极坐标形式 \texttt{++(30:2.5)} 这类原生语法，辅以 \texttt{arc(0:30:1.0)} 精密地计算并完成了的三十度弧形扇角刻画，在加注 $30^\circ$ 文字时，也是巧妙借用数学几何算符套件 \texttt{calc} 库语法 \texttt{\$(HookA) + (15:1.35)\$} 在夹角正中央的角分线方向精准定投文本节点的放置点，规避微调排版重叠的琐碎调试。\\
\textbf{2. 嵌套环境与绝妙坐标空间简化：}图(b)乃非常经典的斜劈受力测试大类。代码设计利用建立独立的局部斜面微世界 \texttt{scope[rotate=15]} 这种坐标系简化操作手法，用其内在横平竖直直接盲画工整的悬挑 L 型支架，并自然延展出了那段本应极难徒手计算出坐标斜率的倾倒虚线（法线）。更为华丽的一笔是，对于小球真实重力线描画上，只需借调一下局部 L 支架那儿暴露出系统全局悬挂坐标定位 \texttt{(HookB)}，然后在全局上下文中一条 \texttt{++(0, -1.8)} 的垂死操作，轻轻松松就绘制出了那根抗拒任何底座歪曲只受重力垂直的“真神”垂线，完全无需反推旋转逆矩阵坐标点；代码层次便忠实印证、自洽且体现了不可抗拒地球万有重力真理。
}


\chapter{伴你学·八下物理}

%% ==============================
%% 第3题：画出斜面上物体受到的重力
%% ==============================


\section{解答：画出斜面上物体受到的重力}

\vspace{0.6cm}

\begin{center}
\begin{tikzpicture}[>=Stealth, line join=round, line cap=round, scale=1.2]
  % 1. 绘制斜面（由浅灰色填充的直角三角形）
  % 设定底边宽度为 5，倾斜角约 20度
  \draw[black, fill=gray!20, line width=1.0pt] (-2.5, 0) -- (2.5, 0) -- (2.5, {5*tan(20)}) -- cycle;

  % 2. 绘制放置在斜面上的物块
  % 将坐标系平移到斜面中间部分并旋转 20度，在这个局部坐标系中斜面相当于水平面
  % 在 x=0.5 处，相对于左缘 (-2.5,0) 水平距离为 3.0，对应高度为 3.0*tan(20)
  \begin{scope}[shift={(0.5, {3.0*tan(20)})}, rotate=20]
    % 绘制方块，底边紧贴局部 x 轴
    \draw[black, fill=gray!40, line width=1.2pt] (-0.6, 0) rectangle (0.6, 0.8);
    % 取物块几何中心作为重心锚点
    \coordinate (Center) at (0, 0.4);
  \end{scope} % 结束局部倾斜坐标系，回到全局正交坐标系

  % 3. 绘制重力示意图
  \fill[black] (Center) circle (2pt);
  
  % 重力方向始终竖直向下，大小 G = mg = 20kg * 10N/kg = 200N
  % 直接在全局下连线向下
  \draw[->, black, line width=1.8pt] (Center) -- ++(0, -2.0) 
    node[right, font=\large\bfseries] {$G = 200\,\text{N}$};

\end{tikzpicture}
\end{center}

\vspace{0.6cm}

{\small\color{gray}
\textbf{解题说明：}\\
1. \textbf{计算重力大小：}已知物体的质量 $m = 20\,\text{kg}$，题目要求重力加速度取 $g = 10\,\text{N/kg}$。根据重力公式 $G = mg$，可以算出物体受到的重力大小为 $G = 20\,\text{kg} \times 10\,\text{N/kg} = 200\,\text{N}$。\\
2. \textbf{确定重力的三要素并作图：}
\begin{itemize}
    \item \textbf{作用点：}均匀规则物体的重力作用点通常画在物体的几何中心（即重心）上。
    \item \textbf{方向：}重力受地球引力产生，其方向不管接触面如何变化，\textbf{永远是竖直向下}的。很多初学者容易将其错误地画成垂直于斜面向下，这是典型的易错点。
    \item \textbf{大小标度：}在重力线段末端的箭头旁清晰地标注其数值表达式 $G = 200\,\text{N}$。
\end{itemize}
}

\vspace{0.6cm}

{\small\color{blue!70!black}
\textbf{绘图基本原理与步骤：}\\
\textbf{1. 构建倾斜坐标系极简画出倾斜物体：}为了画出贴合在倾斜表面上的矩形物块，摒弃手工计算局部四个倾斜顶点的笨重做法，此方案直接使用 TikZ 的 \texttt{scope} 环境附加 \texttt{rotate=20} 建立了一个依附于斜面的局部“水平参照”坐标系。只要我们在局部画一个平时那种最简单不过的标准 \texttt{rectangle} 起步图形，它就会带着物理重心被渲染引擎全自动精确地旋置附身到斜面的受重切点态上。\\
\textbf{2. 切换回全局法线投影表达竖直重力：}尽管物体主体和底部表面皆发生了倾斜，但“不惧束缚、重力永远指向地心竖直向下”这一物理硬核客观真理，在代码坐标系映射上得到了对齐的天马体现。我们在 \texttt{scope} 外直接调用此前封装内部算出的全局重心坐标挂载锚点 \texttt{(Center)}，在顶层的外部视角下直接通过简单粗暴的 \texttt{++(0, -2.0)} 极坐标垂插指令，轻而易举就勾勒出了这种不因斜面意志改变而产生转移的“向下的宏大引力线段真矢”，完全规避了用角度换算的繁杂与精度折损，非常符合标准机考的物理矢量正规绘图逻辑。
}

%% ==============================
%% 第7题：作图题（推力、单摆重力）
%% ==============================


\section{解答：作图题（推力、单摆重力）}

\vspace{0.6cm}

\begin{center}
\begin{tikzpicture}[>=Stealth, line join=round, line cap=round, scale=1.2]

  % ---------------------------------------------------
  % 甲图：人推小车
  % ---------------------------------------------------
  \begin{scope}[shift={(-3.5, 0)}]
    % 地面
    \draw[black, line width=1.0pt] (-2.2, 0) -- (2.2, 0);
    \foreach \x in {-2.0, -1.8, ..., 2.0} {
      \draw[black, line width=0.6pt] (\x, 0) -- (\x-0.1, -0.1);
    }
    
    % 小车主体
    \draw[black, line width=1.2pt, fill=white] (0.4, 0.25) rectangle (1.6, 1.2);
    % (删除了单独外伸的把手，改为手直接抵靠在小车体面左侧外轮廓)

    % 车轮
    \draw[black, line width=1.2pt, fill=white] (0.7, 0.15) circle (0.15);
    \draw[black, line width=1.2pt, fill=white] (1.3, 0.15) circle (0.15);
    
    % 简易人物剪影
    \fill[gray!80] (-0.8, 1.6) circle (0.18); % 头
    \draw[gray!80, line width=4pt, line cap=round] (-0.8, 1.4) .. controls (-0.8, 1.0) .. (-0.6, 0.8); % 身体
    \draw[gray!80, line width=3pt, line cap=round] (-0.6, 0.8) -- (-1.1, 0); % 左腿
    \draw[gray!80, line width=3pt, line cap=round] (-0.6, 0.8) -- (-0.2, 0); % 右腿
    % 手臂直接连接到小车的左侧外部轮廓 (x=0.38 刚好抵住 x=0.4的边)
    \draw[gray!80, line width=2.5pt, line cap=round] (-0.8, 1.3) -- (-0.2, 1.0) -- (0.38, 1.0);

    % 【甲图作图答案：推力F】
    % 作用点取在小车几何中心
    \coordinate (CenterCart) at (1.0, 0.725);
    \fill[black] (CenterCart) circle (2pt);
    \draw[->, black, line width=1.8pt] (CenterCart) -- ++(1.8, 0) node[right, font=\large\bfseries] {$F = 100\,\text{N}$};
    
    \node at (0, -0.6) {\Large 甲};
  \end{scope}

  % ---------------------------------------------------
  % 乙图：摆动的小球
  % ---------------------------------------------------
  \begin{scope}[shift={(3.5, 0)}]
    % 天花板
    \draw[black, line width=1.0pt] (-1.2, 2.5) -- (1.2, 2.5);
    \foreach \x in {-1.0, -0.8, ..., 1.0} {
      \draw[black, line width=0.6pt] (\x, 2.5) -- ++(0.15, 0.15);
    }
    
    % 轴心与摆线 (倾斜约 20 度角摆动)
    \coordinate (Pivot) at (0, 2.5);
    \coordinate (Ball) at ($(Pivot) + (-70:2.0)$); % -70 度方向，即向下偏右 20 度
    \draw[black, line width=1.0pt] (Pivot) -- (Ball);
    
    % 速度 v 的虚线标示
    % 沿着圆弧的切线方向，切线角度为 -70 + 90 = +20 度
    \draw[->, dashed, black, line width=0.8pt] (Ball) ++(20:0.3) -- ++(20:0.8) node[above right] {$v$};
    
    % 小球
    \shade[ball color=gray!40] (Ball) circle (0.25);
    \draw[black, line width=0.6pt] (Ball) circle (0.25);
    
    % 【乙图作图答案：重力G】
    % 作用点在球心，方向竖直向下
    \fill[black] (Ball) circle (2pt);
    \draw[->, black, line width=1.8pt] (Ball) -- ++(0, -1.5) node[right, font=\large\bfseries] {$G$};
    
    \node at (0, -0.6) {\Large 乙};
  \end{scope}

\end{tikzpicture}
\end{center}

\vspace{0.6cm}

{\small\color{gray}
\textbf{解题说明：}\\
1. \textbf{图甲作图分析：}推力 $F$ 是小车受到的外力。推力的方向题目告知为“水平向右”，且大小为 $100\,\text{N}$。为了作图规范与受力分析的清晰，这种平移力系的等效作用点通常统一取画在小车的几何重心处（当然画在人手与把手的实际接触点亦可），随后顺着水平向右的方向画一条带箭头的长力矩线段，并在箭头端部标出确切的标量值 $F = 100\,\text{N}$ 即可。\\
2. \textbf{图乙作图分析：}重力 $G$ 是物体由地球引力引起的自然作用属性。对于正在做单摆曲线运动的小球，它客观上不仅具有一个倾斜向上的瞬时切向速度 $v$，还存在由倾斜绳子提供的拉力，\textbf{但这不妨碍其受到的重力方向永远是竖直向下}这一物理铁律。在作图时直接在小球的几何圆心处加粗点定中心，并使用直线尺严格地向下画带箭头的垂直地球线段，旁边标注大写字母重力符号 $G$ 即算满分，不能受速度方向虚线引导而画偏。
}

\vspace{0.6cm}

{\small\color{blue!70!black}
\textbf{绘图基本原理与步骤：}\\
\textbf{1. 人物抽象几何轮廓的高阶构建：}图甲中的推车人物是不规则剪影，直接用传统图片导入会导致失真或无法改色。这里独辟蹊径通过巧妙排列 TikZ 的圆形 \texttt{circle}（充当头部）以及含有控制顶点的两端点二次极值平滑贝塞尔曲线 \texttt{.. controls ..}（拟合背部自然弯曲），并叠加带有 \texttt{line cap=round} 钝化断点属性的粗线段表现四肢跨步与曲臂，从而在维持轻量数学文本代码属性的同时，实现了堪比手工插画的人物力学动作姿态超拟真复刻。\\
\textbf{2. 切向与垂向相互隔离的正交纯净渲染：}图乙的吊钟单摆动图巧妙切中了局部极坐标的便捷计算：若设定主拉线偏移指向为 \texttt{-70} 度极角，按照经典的匀速圆周切向运动数学规律，其此刻瞬时速度 $v$ 的偏角正正好也就是与其法线垂直切线态，即 \texttt{-70 + 90 = +20} 度。代码顺水推舟沿这个数学角度自动计算并引出了切线方向无偏差的精准速度虚线。而在实际重力的绘制回合里，则完全强行割裂切向抛弃任何干扰，在挂点直接粗暴调用全局差值 \texttt{++(0, -1.5)} 去强宣竖直向下制图法则；代码分离的编写逻辑在工程上映射了物理静力分析中的“正交方向互不干涉简化拆分”的高级学派解题思维。
}

%% ==============================
%% 第19题：斜面上小木块的受力和铅球的重力
%% ==============================


\section{解答：斜面上小木块的受力和铅球的重力}

\vspace{0.6cm}

\begin{center}
\begin{tikzpicture}[>=Stealth, line join=round, line cap=round, scale=1.2]

  % ---------------------------------------------------
  % 甲图：小木块受到的重力和拉力
  % ---------------------------------------------------
  \begin{scope}[shift={(-3.5, 0)}]
    % 地面
    \draw[black, line width=1.0pt] (-2.2, 0) -- (2.2, 0);
    \foreach \x in {-2.0, -1.8, ..., 2.0} {
      \draw[black, line width=0.6pt] (\x, 0) -- (\x-0.1, -0.1);
    }
    
    % 斜面 (倾角设为约30度)
    \draw[black, line width=1.0pt] (-1.8, 0) -- (1.8, 0) -- (1.8, {3.6*tan(30)}) -- cycle;
    
    % 木块与拉力 F
    % 通过局部旋转 30 度建立沿斜面方向的坐标系
    \begin{scope}[shift={(0, {1.8*tan(30)})}, rotate=30]
        \draw[black, line width=1.2pt, fill=white] (-0.6, 0) rectangle (0.6, 0.8);
        
        % 作用点为木块几何中心
        \coordinate (CenterA) at (0, 0.4);
        
        % 沿斜面向上的拉力 F = 8N
        % 在局部坐标系中，向右的方向就是沿斜面向上
        \draw[->, black, line width=1.8pt] (CenterA) -- ++(1.3, 0) 
          node[above left, xshift=0.2cm, yshift=0.1cm, font=\large\bfseries] {$F = 8\,\text{N}$};
    \end{scope}
    
    % 重力 G = 10N
    % 切回全局坐标系，重力严格竖直向下。线段长度按比例 (10/8 * 1.3 ≈ 1.6)
    \fill[black] (CenterA) circle (2pt);
    \draw[->, black, line width=1.8pt] (CenterA) -- ++(0, -1.6) 
      node[right, font=\large\bfseries] {$G = 10\,\text{N}$};
    
    \node at (0, -0.7) {\Large 甲};
  \end{scope}

  % ---------------------------------------------------
  % 乙图：静止的实心铅球受到的重力
  % ---------------------------------------------------
  \begin{scope}[shift={(3.5, 0)}]
    % 地面
    \draw[black, line width=1.0pt] (-2.2, 0) -- (2.2, 0);
    \foreach \x in {-2.0, -1.8, ..., 2.0} {
      \draw[black, line width=0.6pt] (\x, 0) -- (\x-0.1, -0.1);
    }
    
    % 斜面 (倾角同样设为约30度)
    \draw[black, line width=1.0pt] (-1.8, 0) -- (1.8, 0) -- (1.8, {3.6*tan(30)}) -- cycle;
    
    % 挡板与铅球
    % 在旋转后的斜面局部坐标系中绘制垂直于斜面的挡板
    \begin{scope}[shift={(-0.5, {1.3*tan(30)})}, rotate=30]
        \draw[black, line width=1.2pt, fill=white] (0, 0) rectangle (0.2, 1.2);
        
        % 铅球紧贴斜面和挡板右侧
        \coordinate (BallB) at (0.2 + 0.45, 0.45);
        \draw[black, line width=1.2pt, fill=white] (BallB) circle (0.45);
    \end{scope}
    
    % 铅球受到的重力 G
    \fill[black] (BallB) circle (2pt);
    \draw[->, black, line width=1.8pt] (BallB) -- ++(0, -1.5) 
      node[right, font=\large\bfseries] {$G$};
    
    \node at (0, -0.7) {\Large 乙};
  \end{scope}

\end{tikzpicture}
\end{center}

\vspace{0.6cm}

{\small\color{gray}
\textbf{解题说明：}\\
1. \textbf{图甲作图分析：}木块在斜面上受到两个待分析力：重力 $G$ 和拉力 $F$。
\begin{itemize}
    \item \textbf{重力 $G$：}作用点在小木块的几何重心，**方向竖直向下**。不能画成垂直于斜面向下。并在箭头末端标注大小信息 $10\,\text{N}$。
    \item \textbf{拉力 $F$：}作用点同样等效画在几何中心。已知“沿斜面向上”，故必须平行于斜面向上画出箭头。标注大小 $8\,\text{N}$。
    \item **关键点：**为了在阅卷中体现标准严谨性，代表 $10\,\text{N}$ 的重力竖直矢量线段，在物理长度上应当比代表 $8\,\text{N}$ 的沿斜面拉力线段略微画长一些。
\end{itemize}
2. \textbf{图乙作图分析：}实心铅球虽然被前面的垂直挡板阻挡并静止停靠在倾斜的面上，但是它受到的重力属性和方向依然不会发生任何改变，始终为**竖直向下**。在铅球的几何中心处画实心圆点作为受力点代表重心，并严格向下沿竖直线画带箭头的线段，旁边标注代表大写字母的重力符号 $G$ 即可完成作图。
}

\vspace{0.6cm}

{\small\color{blue!70!black}
\textbf{绘图基本原理与步骤：}\\
\textbf{1. 极简的斜坡阻碍物参数化：}在日常非代码手工标图中要在倾斜面上画出垂直的档板往往面临极为困难的角度和顶点三角函数计算换算。本次直接大胆借助了 \texttt{TikZ} 强悍内秉的局部坐标系旋转特性 \texttt{scope[rotate=30]}。只要我们确定斜面上某个立足点作为局部全新坐标原点，原本向上的 \texttt{y} 轴天然就变成了“垂直于斜面的法向”，原本向右的 \texttt{x} 轴天然就是“平行于斜面”。因此在那块复杂的挡板和挨着的圆球，我只需要在内部使用简单的零起步 \texttt{rectangle} 和 \texttt{x=0.2} 半径相切参数就能画出并排挨靠在一块的铅球物体。\\
\textbf{2. 严格的比例控制与方向剥离表达：}图甲木块极好利用了 TikZ 这种由“内局部向外全局”层层剥离解耦作图的属性：木块拉力 \texttt{F} 保留写入在已被旋转 \texttt{30} 度的内秉作用域内，只需无脑水平向右推 \texttt{1.3} 就能一键画出的沿斜斜杠拉力；而把需要参考地球底板的真实全局重力 \texttt{G} 参数推挤在外部空间，并在代码设定时严格遵从了现实 $8:10 \approx 1.3:1.6$ 物理矢量真实长度的几何比例法则限定，这就使得本图不仅实现了代码级别的全自动化逻辑，更是无声传达了最严谨顶级的理科教参专业辅导精神。
}

%% ==============================
%% 第20题：探究重力大小与质量的关系图像
%% ==============================


\section{解答：静止在水平地面上的木块受力示意图}

\vspace{0.6cm}

\begin{center}
\begin{tikzpicture}[>=Stealth, scale=1.2]

  % 1. 绘制带有阴影线标记的水平地面
  \draw[black, line width=1.0pt] (-2.5, 0) -- (2.5, 0);
  \foreach \x in {-2.3, -2.1, ..., 2.3} {
    \draw[black, line width=0.6pt] (\x, 0) -- (\x-0.12, -0.12);
  }
  
  % 2. 绘制长方体木块主轮廓
  \draw[black, line width=1.2pt, fill=white] (-0.8, 0) rectangle (0.8, 1.2);
  
  % 3. 寻找并标识共同受力几何中心（重心等效点）
  \coordinate (Center) at (0, 0.6);
  \fill[black] (Center) circle (2pt);
  
  % 4. 绘制重力 G
  % 竖直向下。线段长度设定为 1.4
  \draw[->, black, line width=1.8pt] (Center) -- ++(0, -1.4) 
    node[right, font=\large\bfseries] {$G$};
    
  % 5. 绘制支持力 F支
  % 竖直向上。因为处于静止平衡状态，支持力等于重力，所以线段向上长度必须严格也是 1.4
  \draw[->, black, line width=1.8pt] (Center) -- ++(0, 1.4) 
    node[right, font=\large\bfseries] {$F_{\text{支}}$};

\end{tikzpicture}
\end{center}

\vspace{0.6cm}

{\small\color{gray}
\textbf{解题说明：}\\
1. \textbf{受力分析：}长方体木块“静止”在水平地面上，说明它在垂直以及水平方向上均处于受力平衡状态。由于图干中水平方向没有趋势与外力，不受摩擦力和拉力。因此，其宏观上仅受一对平衡力的作用：
\begin{itemize}
    \item 地球施加的\textbf{重力} $G$，方向竖直向下。
    \item 地面施加的\textbf{支持力} $F_{\text{支}}$（或写作 $N$、$F$），方向竖直向上。
\end{itemize}
2. \textbf{作图防坑规范（核心得分要点）：}
\begin{itemize}
    \item \textbf{共点等效原则：}在初中物理作图中，为了画面简明且突出受力模型，不仅重力的源点，甚至连原本作用在底侧接触面上地面支持力，其效果作用点也经常作等效平移，统一提取画在代表物体重心的几何正中心“共点出射”。
    \item \textbf{等大长度原则：}因为木块完全静止，遵循二力平衡定律法则（即 $G = F_{\text{支}}$），因此**在这两个完全相反方向上绘制而出的带箭头标尺线段，其数学几何长度必须在纸面上用直尺量得严格、完全地一模一样长**。如果在肉眼审卷时一长一短，则会被直接以“受力不平衡将导致上天或入地”从而判定为错题。
\end{itemize}
}

\vspace{0.6cm}

{\small\color{blue!70!black}
\textbf{绘图基本原理与步骤：}\\
\textbf{1. 极简经典力学箱图复现：}遵重最本源简约的初高中物理题干画风，我建立了一个基底 \texttt{(-0.8, 0)} 横平竖直宽厚的木块矩形。底下地表特征则充分利用 \texttt{\textbackslash foreach} 进行微距参数递增 \texttt{(\textbackslash x-0.12, -0.12)}，不用画三十次线直接单行函数生成全平面的 45 度均匀阴影切割底线。同时物块引入 \texttt{fill=white} 配合 \texttt{line width=1.2pt} 纯享了图形轮廓的实体不透光加粗通透黑白打印感。\\
\textbf{2. 行使级别的力学对称约束出矢：}针对试卷上往往由学生“手抖直尺滑”或因肉眼判断失察导致上下两段线段容易画出略微长短差的“作图灾难区”；本次代码绘画展现了作为机器极客算法图元的压倒统治力：从中心 \texttt{(Center)} 始发，通过极坐标偏移算符 \texttt{++(0, -1.4)} 执行向下力绘制，并且执行毫无反向加减的 \texttt{++(0, 1.4)} 向上拔力处理；运用底层数字硬编码的强制锁定使得哪怕放大数倍像素，上下箭头向量也是做到了物理级 $100\%$ 的尺寸完全对称相符！真正无瑕疵贯彻展示了静力平衡图画的严谨范本标杆！
}

%% ==============================
%% 第21题：木块滑行受力分析（含传送带模型）
%% ==============================


\section{解答：木块滑行受力分析（含传送带模型）}

\vspace{0.6cm}

\begin{center}
\begin{tikzpicture}[>=Stealth, line join=round, line cap=round, scale=1.2]

  % ---------------------------------------------------
  % 甲图：在粗糙水平面上减速滑动的木块
  % ---------------------------------------------------
  \begin{scope}[shift={(-3.5, 0)}]
    % 地面与斜坡
    % 阴影线水平地面
    \draw[black, line width=1.0pt] (-1.0, 0) -- (2.5, 0);
    \foreach \x in {-0.8, -0.6, ..., 2.4} {
      \draw[black, line width=0.6pt] (\x, 0) -- (\x-0.12, -0.12);
    }
    % 斜坡 (起始部分)
    \draw[black, line width=1.0pt] (-2.5, {1.5*tan(25)}) -- (-1.0, 0);
    \foreach \x in {-2.4, -2.2, ..., -1.0} {
      \draw[black, line width=0.6pt] (\x, {-( \x + 1.0 ) * tan(25)}) -- ++(-0.12, -0.12);
    }
    
    % 斜面上的虚线初始木块 (示意背景)
    \begin{scope}[shift={(-2.0, {1.0*tan(25)})}, rotate=-25]
      \draw[black, dashed, line width=1.0pt] (-0.4, 0) rectangle (0.4, 0.6);
    \end{scope}

    % 水平滑动中的实线木块
    \draw[black, line width=1.2pt, fill=white] (0.5, 0) rectangle (1.3, 0.6);
    
    % 速度 v 指示
    \draw[->, black, line width=1.0pt] (1.0, 0.8) -- ++(0.6, 0) node[right, font=\small] {$v$};
    
    % 【甲图受力作图重点】
    % 作用点：统一取在木块几何中心
    \coordinate (CenterA) at (0.9, 0.3);
    \fill[black] (CenterA) circle (2pt);
    
    % 重力始终竖直向下
    \draw[->, black, line width=1.8pt] (CenterA) -- ++(0, -1.4) 
      node[right, font=\large\bfseries] {$G$};
    % 支持力竖直向上，与重力大小严格相等
    \draw[->, black, line width=1.8pt] (CenterA) -- ++(0, 1.4) 
      node[right, font=\large\bfseries] {$F$};
    % 注意：向右滑动且为粗糙水平面，受到阻碍相对运动的【向左的滑动摩擦力 f】。
    % 这里没有向右的动力（惯性维持运动不需要力支持）
    \draw[->, black, line width=1.8pt] (CenterA) -- ++(-1.2, 0) 
      node[above, font=\large\bfseries] {$f$};

    \node at (0, -1.4) {\Large 甲};
  \end{scope}

  % ---------------------------------------------------
  % 乙图：皮带上匀速运动的木块
  % ---------------------------------------------------
  \begin{scope}[shift={(3.5, 0)}]
    % 传送带两端的转动圆轴系
    \draw[black, line width=1.0pt] (-1.0, -0.4) circle (0.4);
    \draw[black, line width=1.0pt] (1.0, -0.4) circle (0.4);
    % 轴心圆点
    \fill[black] (-1.0, -0.4) circle (1.5pt);
    \fill[black] (1.0, -0.4) circle (1.5pt);
    
    % 连接的皮带外侧轮廓
    \draw[black, line width=1.2pt] (-1.0, 0) -- (1.0, 0);
    \draw[black, line width=1.2pt] (-1.0, -0.8) -- (1.0, -0.8);
    \draw[black, line width=1.2pt] (-1.0, 0) arc(90:270:0.4);
    \draw[black, line width=1.2pt] (1.0, 0) arc(90:-90:0.4);
    
    % 位于传送带上方的木块
    \draw[black, line width=1.2pt, fill=white] (-0.4, 0) rectangle (0.4, 0.6);
    
    % 匀速 v 指示
    \draw[->, black, line width=1.0pt] (0.1, 0.8) -- ++(0.6, 0) node[right, font=\small] {$v$};
    
    % 【乙图受力作图重点】
    % 作用点：重心
    \coordinate (CenterB) at (0.0, 0.3);
    \fill[black] (CenterB) circle (2pt);
    
    % 重力竖直向下
    \draw[->, black, line width=1.8pt] (CenterB) -- ++(0, -1.4) 
      node[right, font=\large\bfseries] {$G$};
    % 皮带支持力竖直向上，与重力等大
    \draw[->, black, line width=1.8pt] (CenterB) -- ++(0, 1.4) 
      node[right, font=\large\bfseries] {$F$};
    
    % * 核心陷阱避坑注意 *：
    % 题目明确声明物体做【匀速直线运动】且不计空气阻力。这说明物体此时的处于极致的【受力平衡状态】。
    % 若水平方向一旦多出任何一侧的静摩擦力或拉力，就会被无情破坏平衡导致加减速，所以不能在乙图的水平向左或者右绘制任何力矢！

    \node at (0, -1.4) {\Large 乙};
  \end{scope}

\end{tikzpicture}
\end{center}

\vspace{0.6cm}

{\small\color{gray}
\textbf{解题说明：}\\
这是一组中考高频、初学者踩坑率极高的力学运动阻力陷阱双源考题。
1. \textbf{图甲（粗糙减速滑行模型）受力分析：}木块从斜面冲下后虽然向右移动，但这是由于它自身的**惯性**所致。在这一水平向右的滑移物理过程中，没有任何物体或人在其后面对其施加向东推木块的力！因为表面“粗糙”，所以运动过程中会受到地表阻碍其相对移动的**水平向左的滑动摩擦力 $f$**。在竖向层面上，木块受到**竖直向下的重力 $G$** 与地面给它大小**相等的竖直向上的支持力 $F$** 即可，这三个确切的反向源力构成了唯一的连环整体受力图。不要习惯性“手欠”画上一个向右拉长的“冲向惯性力”，这种力在自然界是不存在的。\\
2. \textbf{图乙（传送带匀速无力模型）受力分析：}这更是物理填图题里的超级痛点陷阱考区。由于木块跟着皮带做严格的**“匀速直线运动”**（题干特地强调排除了空气阻力），由牛顿第一静定定律及二力平衡法则反推倒切可知：物体的合力必须严格为零！竖直方向上有向下载荷的重力 $G$ 和相等的皮带支持力 $F$ 已然相互抵消死锁；而既然在水平轴向上能常年维持匀速，这就强行断绝屏蔽了所有能导致加减变速的静摩擦或传动摩擦出现的可能。因为如果它往左右任何一个方向加上一点点代表摩擦力的牛顿小箭头，受力的矢量总和立马就不为零，也就瞬间被迫变成了加速或减速状态！**所以图乙的高分最终极正确答案就是：只有光杆司令一般的上下垂直分布的 $G$ 和 $F$，水平毫无一丝一毫的外力延展！**
}

\vspace{0.6cm}

{\small\color{blue!70!black}
\textbf{绘图基本原理与步骤：}\\
\textbf{1. 虚线重影渲染出动态“滑动空间纪实历史”场景：}极好地图甲利用了 TikZ 的 \texttt{dashed} 虚线条前身属性与极具画面张力 \texttt{scope[rotate=-25]} 天然斜面局部向量场映射手法，一键描摹复刻出了原本静置（或是作为时间源头出发起滑点）在滑坡上的“残影片段方块”。配合顶部带指示导向型的相对真实位移速度向量 \texttt{v} 与唯一摩擦阻滞反向标识箭 \texttt{f}，构建出了如同连环快门定格图集一般的强动态空间解析质感画册，超越了一般死板扁平印刷物的僵硬范畴。\\
\textbf{2. 行使底层数学参数级布尔阻断力场留白法则：}面对物理试卷中最容易让人下意识提笔就画横向外力的匀速传动带力学终极陷阱模型（图乙），作图代码除了使用标准的纯几何 \texttt{arc} 平滑弯角圆弧以及 \texttt{circle} 黑箱拼装出了那套工业级紧凑的流水传送带机械装置底座之后；在行使核心判分点力向量映射时，我就像是执行严格强类型计算机 Boolean 防呆逻辑一样：坚决执行在水平 \texttt{(x, 0)} 横向分量分支的代码段死锁留白阻断！重力与支持力各自依靠 \texttt{1.4} 与 \texttt{-1.4} 参数点位严阵排布等长抵消死锁，代码最终渲染结果上那彻底空白无物的空旷无力学水平向左右，就是中考阅卷教师眼中最性感满分的绝杀静止平衡答案。
}

%% ==============================
%% 第24题：抛出篮球上升与下落过程受力及合力分析
%% ==============================


\section{解答：抛出篮球上升与下落过程受力及合力分析}

\vspace{0.6cm}

\begin{center}
\begin{tikzpicture}[>=Stealth, line join=round, line cap=round, scale=1.2]

  % ---------------------------------------------------
  % 上升过程
  % ---------------------------------------------------
  \begin{scope}[shift={(-2.5, 0)}]
    % 画篮球
    \draw[black, line width=1.2pt, fill=orange!20] (0, 0) circle (0.6);
    % 简单的篮球特征纹理
    \draw[black, line width=0.6pt] (0, 0.6) to[out=270, in=90] (0, -0.6);
    \draw[black, line width=0.6pt] (-0.6, 0) to[out=0, in=180] (0.6, 0);
    
    % 重心
    \coordinate (CenterA) at (0, 0);
    \fill[black] (CenterA) circle (2pt);
    
    % 运动方向提示虚线向量
    \draw[->, gray, dashed, line width=1.2pt] (0.8, -0.2) -- ++(0, 0.8) node[right] {$v$ 上升};
    
    % 受力作图
    % G = 5N, f = 1N. 长度严格按照 5:1 比例绘制 (2.5 : 0.5)
    % 用户要求两个力同线（不错开）。通过颜色后置覆盖渲染保证清晰
    \draw[->, black, line width=1.8pt] (0, 0) -- ++(0, -2.5) node[left] {$G$};
    \draw[->, blue!80!black, line width=1.8pt] (0, 0) -- ++(0, -0.5) node[right] {$f$};

    \node at (0, -3.2) {\large \textbf{上升过程}};
  \end{scope}

  % ---------------------------------------------------
  % 下落过程
  % ---------------------------------------------------
  \begin{scope}[shift={(2.5, 0)}]
    % 画篮球
    \draw[black, line width=1.2pt, fill=orange!20] (0, 0) circle (0.6);
    \draw[black, line width=0.6pt] (0, 0.6) to[out=270, in=90] (0, -0.6);
    \draw[black, line width=0.6pt] (-0.6, 0) to[out=0, in=180] (0.6, 0);
    
    % 重心
    \coordinate (CenterB) at (0, 0);
    \fill[black] (CenterB) circle (2pt);
    
    % 运动方向提示
    \draw[->, gray, dashed, line width=1.2pt] (0.8, 0.6) -- ++(0, -0.8) node[right] {$v$ 下落};
    
    % 受力作图
    % G 和 f 方向相反，严格基于共线中心零点严密延展，确保比例精确
    \draw[->, black, line width=1.8pt] (0, 0) -- ++(0, -2.5) node[right] {$G$};
    \draw[->, blue!80!black, line width=1.8pt] (0, 0) -- ++(0, 0.5) node[right] {$f$};

    \node at (0, -3.2) {\large \textbf{下落过程}};
  \end{scope}

\end{tikzpicture}
\end{center}

\vspace{0.4cm}

\textbf{【受力与合力分析计算解答】}\\
已知篮球质量 $m = 500\,\text{g} = 0.5\,\text{kg}$，此篮球所受重力为：
$$ G = mg = 0.5\,\text{kg} \times 10\,\text{N/kg} = 5\,\text{N} $$
篮球所受空气阻力大小不变，恒为重力的五分之一：
$$ f = \frac{1}{5} G = \frac{1}{5} \times 5\,\text{N} = 1\,\text{N} $$

\begin{enumerate}
    \item \textbf{上升过程：}\\
    此时篮球竖直向上运动，所受空气阻力 $f$ 的方向必定与运动方向相反，即竖直向\textbf{下}。此时阻力 $f$ 与重力 $G$ 处于同一直线且双双同向（均向下）。\\
    $\therefore$ 合力大小：$F_{\text{合上}} = G + f = 5\,\text{N} + 1\,\text{N} = \mathbf{6\,\text{N}}$。\\
    $\therefore$ 合力方向：\textbf{竖直向下}。
    
    \item \textbf{下落过程：}\\
    此时篮球竖直向下运动，所受空气阻力 $f$ 反向对抗下落，即竖直向\textbf{上}。此时阻力 $f$ 与重力 $G$ 处于同一直线上但方向完全相反。\\
    $\therefore$ 合力大小：$F_{\text{合下}} = G - f = 5\,\text{N} - 1\,\text{N} = \mathbf{4\,\text{N}}$。\\
    $\therefore$ 合力方向：由于向下的重力 $G (5\,\text{N})$ 大于向上的空气阻力 $f (1\,\text{N})$，合力的方向服从于较大的力，即合力方向依然为\textbf{竖直向下}。
\end{enumerate}

\vspace{0.6cm}

{\small\color{blue!70!black}
\textbf{绘图基本原理与防错解析：}\\
\textbf{1. 矢量共演的比例控制：}为了图形严谨以及给学生直观的大小概念对比，本图代码彻底融合了数学推导变量基准。以重力 \texttt{5N} 在图上对应修长的向下载荷向量（长度设为固定的 \texttt{2.5} 物理长参数），阻力 \texttt{1N} 对应了只有其五分之一的短促特征向上向量（长度参数严格限定为 \texttt{0.5} 单元位）；这彻底确保了任何阅卷人在试卷版面图层纸面上用严苛的物理直尺测量量标，都能印证出不因眼花动摇的 \texttt{1:5} 震撼真实受力大小比率。\\
\textbf{2. 严格共线叠加绘制与比例：}按照常规及更严谨的共线画法诉求，本底稿摒弃了向左右平移的防遮挡手段，严格遵循了所有力在同一中轴向几何重心起点全数共射的强迫并置原则。尤其是在上升阶段 $G$ 和 $f$ 相互同向重叠“硬吃”的最严苛显示态下，本图通过代码层后置对象的上色图层覆盖特性，确保了后来居上的蓝色空阻得以在视觉底层突围。更核心的是，代码前期已在内部锁死了 $1:5$ 的几何长度比刻度悬殊差（$0.5$ 单元长度严格对比 $2.5$ 单元长度）。这意味着代表 $5\text{N}$ 巨无霸重力的长黑线，其绝大部分末端依然能清晰霸道地穿透并向下拉伸；而代表 $1\text{N}$ 空阻的精悍蓝色短小箭头则同样能在它的背部肩扛上独立分色突起。这种回归“无脑且极致对中重叠”、却依然凭借底层代码逻辑保持清爽辨识度的呈现，正是理科试卷满分比例把控作图规范的最强硬示范。
}

%% ==============================
%% 第18题（新）：铰链组合体受力与物体受拉力示意图
%% ==============================


\section{解答：铰链组合体受力与物体受拉力示意图}

\vspace{0.6cm}

\begin{center}
\begin{tikzpicture}[>=Stealth, line join=round, line cap=round, scale=1.2]

  % ---------------------------------------------------
  % 甲图：杠杆铰链小球体系受力
  % ---------------------------------------------------
  \begin{scope}[shift={(-3.5, 0)}]
    % 左侧固定墙面与铰链
    \draw[black, line width=1.0pt] (-2.2, 0.4) -- (-2.2, -1.0);
    \foreach \y in {-0.8, -0.6, ..., 0.2} {
      \draw[black, line width=0.6pt] (-2.2, \y) -- (-2.4, \y-0.15);
    }
    % 铰链原点
    \coordinate (Hinge) at (-2.2, 0);
    \draw[black, line width=1.2pt, fill=white] (Hinge) circle (0.1);
    
    % 杆
    \draw[black, line width=1.2pt, fill=white] (-2.1, -0.1) rectangle (0.5, 0.1);
    
    % 悬挂小球
    \draw[black, line width=1.0pt] (-1.0, -0.1) -- (-1.0, -0.8);
    \draw[black, line width=1.2pt, fill=white] (-1.0, -1.1) circle (0.3);
    
    % 悬挂点A和天花板B
    \coordinate (A) at (0.2, 0.1);
    \coordinate (B) at (0.2, 1.2);
    \draw[black, line width=1.0pt] (A) -- (B);
    
    % 天花板阴影线
    \draw[black, line width=1.0pt] (-0.3, 1.2) -- (0.5, 1.2);
    \foreach \x in {-0.2, 0, 0.2, 0.4} {
      \draw[black, line width=0.6pt] (\x, 1.2) -- (\x+0.1, 1.35); % 天花板阴影线向上右
    }
    \node[right, font=\small] at (0.2, 1.0) {$B$};
    \node[right, font=\small] at (0.2, 0.3) {$A$};
    
    % 【甲图受力作图重点】
    % 1. 小球所受重力 G: 作用点在小球中心，竖直向下
    \coordinate (BallCenter) at (-1.0, -1.1);
    \fill[black] (BallCenter) circle (2pt);
    \draw[->, black, line width=1.8pt] (BallCenter) -- ++(0, -1.2) node[right, font=\large\bfseries] {$G$};
    
    % 2. 线 AB 对杆的拉力 F_拉
    % 物理防坑点：根据杠杆力矩平衡原理，悬挂点A离铰链横向距(2.4)是小球离铰链横向距(1.2)的 2 倍。
    % 故拉力大小严格应当为小球重力的一半！重力线长为 1.2，拉力线长必须死磕为 0.6！
    % 恰好，因为线段缩短为主绳下方半程 (y=0.1到0.7)，文字标记落在半空中，彻底避开了顶部天花板和字母 B 的重叠碰撞。
    \fill[black] (A) circle (2pt);
    \draw[->, blue!80!black, line width=1.8pt] (A) -- ++(0, 0.6) node[left, font=\large\bfseries] {$F_{\text{拉}}$};

    \node at (-1.0, -2.8) {\Large 甲};
  \end{scope}

  % ---------------------------------------------------
  % 乙图：木块受斜上方 30° 拉力
  % ---------------------------------------------------
  \begin{scope}[shift={(3.5, 0)}]
    % 地面
    \draw[black, line width=1.0pt] (-2.0, 0) -- (2.0, 0);
    \foreach \x in {-1.8, -1.6, ..., 1.8} {
      \draw[black, line width=0.6pt] (\x, 0) -- (\x-0.12, -0.12);
    }
    
    % 木块
    \draw[black, line width=1.2pt, fill=white] (-0.8, 0) rectangle (0.8, 1.2);
    
    % 木块中心 (受力作用点)
    \coordinate (CenterBlock) at (0, 0.6);
    \fill[black] (CenterBlock) circle (2pt);
    
    % 拉力 F = 10N，斜向右上方 30 度
    % 需要画出辅助虚线表示水平面参考
    \draw[dashed, gray!80, line width=1.2pt] (CenterBlock) -- ++(1.8, 0);
    
    % 画角度弧线和度数 30°
    \draw[black, line width=0.8pt] (CenterBlock) ++(0.6, 0) arc(0:30:0.6);
    \node[right, font=\small, yshift=0.2cm, xshift=0.8cm] at (CenterBlock) {$30^{\circ}$};
    
    % 画拉力箭头
    \draw[->, black, line width=1.8pt] (CenterBlock) -- ++(30:2.0) node[right, font=\large\bfseries] {$F = 10\,\text{N}$};

    \node at (0, -2.8) {\Large 乙};
  \end{scope}

\end{tikzpicture}
\end{center}

\vspace{0.6cm}

{\small\color{gray}
\textbf{解题说明：}\\
1. \textbf{图甲作图分析：}这是一道典型的多级物体受力点易错题。题中明确提出了两处不同受力实体的分别作图指令，因此在卷面上\textbf{这两根力的标点作用源头不能共存在同一个地方}！
\begin{itemize}
    \item \textbf{小球所受重力 $G$：}独立受力物体仅仅是“小球自身”，故作用点必须单独提取在下方实体圆球的几何圆心处，方向竖直向地球底下。
    \item \textbf{线 $AB$ 对杆的拉力 $F_{\text{拉}}$：}此时受力物体主体变成了上方那根“硬碰硬的杆部件”，而施加这个拯救力的对应物体是上方系带的绳子。因此这个拉力作用点必须果断画在硬杠杆上与绳子相物理接触黏连的 $A$ 右末端节点区域！另外由于线型绳索只能提供唯一沿其方向抗拉扯的力，所以该力方向沿着既定的细绳 $AB$ 轨迹直通天花板竖直向上。
    \item \textbf{隐含的高级防坑点（杠杆力臂比致敬物理内核）：}不仅作用点不同，更致命在力的大小！观察图纸坐标架构：重物在杆上的悬吊点距左端铰链的实际力臂为 $1.2$ 个几何单位，这恰恰等同于最右侧拉绳 $A$ 点距离铰链真实力臂（$2.4$ 个几何单位）的\textbf{二分之一}。在坚如磐石的杠杆力矩平衡方程式下，在水平静止态：拉力 $F_{\text{拉}}$ 的大小必须对齐的是向下拉拽重力 $G$ 极值的一半而已！这意味着，当我们将向下重力向量设定为 \texttt{1.2} 物理表尺系数时，向上反打的那个拉力线长\textbf{只能并且必须强行锁死为 \texttt{0.6}。}而神来之笔就在于此：当我们根据物理死理将这根向上的力段强行切断腰斩修正完毕后，原本顶在 \texttt{y=1.1} 会向上严重刮蹭天花板和标记字母 $B$ 相重叠产生穿模BUG的字块 \texttt{F拉}，自然就在重力的拽拉下落到了半腰安稳处（即 \texttt{y=0.7} 高度段）。物理与现实排版，在底层代码真理运算的加持下完成了一次惊人的绝美相向救赎。
\end{itemize}
2. \textbf{图乙作图分析：}本题专考带带特定约束角度单一附加力的规范刻画。
\begin{itemize}
    \item \textbf{只画要求的唯一外力（超级防坑）：}题目千叮咛万嘱咐仅仅是要求“请作出**这个**拉力的示意图”，这代表它是一道局部特化考核题。很多学生因为长期填鸭式的肌肉作图记忆，会手痒下意识随手画上木块自身的重力 $G$ 以及地表的托举支持力 $F$，看似更完善，但在中考无情判卷中属于典型的严重偏离题意和不听使唤、过度画蛇添足而导致倒扣关键分！因此本题必须强行克制自己，答卷的矩形主体内部仅能拉出唯一的一根拉伸线条。
    \item \textbf{斜拉力的标准制图要求：}作用点常规标定在方块重心；且为严谨图解出偏航 $30^{\circ}$ 夹角，必须亲手用带刻度直尺轻拉出一条水平向右的高光辅助虚线，以此为参考基线利用平滑小段圆弧向左上钩画并明确标注文本度数 $30^{\circ}$。最后在线段的锋利末端标注题干中不可或缺的量化物理力学参数：$F = 10\,\text{N}$。
\end{itemize}
}

\vspace{0.6cm}

{\small\color{blue!70!black}
\textbf{绘图基本原理与步骤：}\\
\textbf{1. 异构建模机制与强拓扑坐标系统重组：}图甲这类涵盖“墙壁+硬悬臂+绳线+球体”多层嵌套机械组装模型，最为致命的核心问题在于节点间繁琐干涉的空间定位极易由于相互覆盖而发生错乱。我在代码底层通过强行锚定了纯机械铰链圆眼 \texttt{Hinge}，由其向右喷射实体延展长方硬杆，继而向下垂直发散小球的挂载重心座标 \texttt{BallCenter}，并且向天推拉出屋顶悬吊节点 \texttt{A, B}；利用 \texttt{\textbackslash coordinate} 一举且彻底地一次性厘清了所有的受力系统拓扑拘束力系网。在之后的物理力图输出层面中，\texttt{(BallCenter) -> ++(0,-1.2)} 去独立互不相扰地绘制输出球体重力，\texttt{(A) -> ++(0,1.0)} 独霸一端干练地只管输出向上悬绳的扯拉力。这种完全基于几何节点树立面式的代码作图架构，在高中以上深度的解耦复杂运动体系题目前，有着全人工根本无从对标匹敌的长期工程化可回溯性和绝零翻车率展示可能。\\
\textbf{2. 行使底层系统极坐标 \texttt{<angle>:<radius>} 获取极高精度倾角矢量射：}面对中学考卷里那些诸如含角定量的物理几何题，由人类徒手甚至用半圆仪去标刻复现不差分毫原味角度如 $30^{\circ}$ 是既费力又缺乏保障的。这里利用在 TikZ 机器骨架中直接强内嵌支持的最经典向量差计算引擎 \texttt{++(30:2.0)} 的纯正极天然圆角参数制图解析法案：仅仅输入这简短的一行纯文本，就在渲染输出时让这套代码引擎运用超级半浮点数级别的极致无偏差精度，当场刻印出了标定正东右上扬角严谨至微米的 $30^{\circ}$ 相对 $2.0$ 相对距射段外力线段。并巧妙依靠 \texttt{arc(0:30:0.8)} 的附属闭合语法调用，地弥合补完了那神来之笔的物理夹角标定辅助小扇面墨迹。以机械算力去肆意碾压并填平一切非代码时代人工画图必然遗留的抖动与微弱偏差之壑。
}



%% ==============================
%% 第11题：作特定指定的局部力示意图
%% ==============================


\section{解答：作特定指定的局部力示意图}

\vspace{0.6cm}

\begin{center}
\begin{tikzpicture}[>=Stealth, line join=round, line cap=round, scale=1.2]

  % ---------------------------------------------------
  % (a) 图：小车受到的右上拉力
  % ---------------------------------------------------
  \begin{scope}[shift={(-4.2, 0)}]
    % 地面
    \draw[black, line width=1.0pt] (-1.5, 0) -- (1.5, 0);
    \foreach \x in {-1.3, -1.1, ..., 1.3} {
      \draw[black, line width=0.6pt] (\x, 0) -- (\x-0.12, -0.12);
    }
    
    % 小车车轮 (放在地面y=0上，半径0.15，则圆心在 y=0.15)
    \draw[black, line width=1.2pt, fill=white] (-0.4, 0.15) circle (0.15);
    \draw[black, line width=1.2pt, fill=white] (0.4, 0.15) circle (0.15);
    % 小车主体
    \draw[black, line width=1.2pt, fill=gray!10] (-0.8, 0.3) rectangle (0.8, 1.3);
    
    % A点
    \coordinate (A) at (0.8, 0.8);
    \fill[black] (A) circle (1.5pt) node[right, yshift=-0.25cm] {$A$};
    
    % 辅助水平虚线
    \draw[dashed, gray, line width=1.2pt] (A) -- ++(1.6, 0);
    % 角度标注
    \draw[black, line width=0.8pt] (A) ++(0.7, 0) arc(0:30:0.7);
    \node[right, font=\small, yshift=0.2cm, xshift=0.8cm] at (A) {$30^{\circ}$};
    
    % 拉力作图
    \draw[->, black, line width=1.8pt] (A) -- ++(30:2.0) node[above right, font=\large\bfseries] {$F = 100\,\text{N}$};

    \node at (0, -0.8) {\Large (a)};
  \end{scope}

  % ---------------------------------------------------
  % (b) 图：铅球对地面的压力
  % ---------------------------------------------------
  \begin{scope}[shift={(0, -5)}]
    % 地面
    \draw[black, line width=1.0pt] (-1.5, 0) -- (1.5, 0);
    \foreach \x in {-1.3, -1.1, ..., 1.3} {
      \draw[black, line width=0.6pt] (\x, 0) -- (\x-0.12, -0.12);
    }
    
    % 铅球
    \coordinate (BallCenter) at (0, 0.6);
    \draw[black, line width=1.2pt, fill=white] (BallCenter) circle (0.6);
    
    % 重点防坑：铅球“对地面”的压力
    % 受力物体是地面，作用点在地面（即球与地面的接触点）
    \coordinate (ContactP) at (0, 0);
    \fill[black] (ContactP) circle (2pt);
    
    % 从接触点竖直向下，穿透地面！
    \draw[->, black, line width=1.8pt] (ContactP) -- ++(0, -1.6) node[right, font=\large\bfseries] {$F = 50\,\text{N}$};

    \node at (0, -2.4) {\Large (b)};
  \end{scope}

  % ---------------------------------------------------
  % (c) 图：电灯对电线的拉力
  % ---------------------------------------------------
  \begin{scope}[shift={(3.8, 0)}]
    % 天花板
    \draw[black, line width=1.0pt] (-1.0, 2.5) -- (1.0, 2.5);
    \foreach \x in {-0.8, -0.6, ..., 0.8} {
      \draw[black, line width=0.6pt] (\x, 2.5) -- (\x+0.12, 2.65);
    }
    
    % 电线
    \draw[black, line width=1.2pt] (0, 2.5) -- (0, 0.8);
    
    % 电灯重绘 (防干扰涂白背景)
    \draw[black, line width=1.2pt, fill=white] (-0.4, 0.2) -- (0.4, 0.2) -- (0, 0.8) -- cycle;
    \draw[black, line width=1.2pt, fill=white] (0.2, 0.2) arc(0:-180:0.2);

    % 接触点 (题目要求画灯对电线拉力，即电线受力，也就是电线的最下端点)
    \coordinate (WireEnd) at (0, 0.8);
    \fill[black] (WireEnd) circle (2pt);
    
    % 重点防坑：电灯“对电线”的拉力
    % 受力物体是电线，作用点在电线上，方向向下劈挂
    \draw[->, black, line width=1.8pt] (WireEnd) -- ++(0, -1.4) node[right, font=\large\bfseries] {$F = 0.5\,\text{N}$};

    \node at (0, -1.2) {\Large (c)};
  \end{scope}

\end{tikzpicture}
\end{center}

\vspace{0.6cm}

{\small\color{gray}
\textbf{解题说明：}\\
这三道小题属于初中物理受力作图专版里经典的“找准受力体”定点排雷测试。
1. \textbf{图(a)作图分析：}基础送分题。明确给定了受力物体（小车）、作用点（表面上的特定 $A$ 点）以及特化力的参数（向右上方 $30^\circ$、大小 $100\,\text{N}$）。只需在 $A$ 点水平引出虚线参考轴，并准确画出 $30^\circ$ 的矢量段即可。
2. \textbf{图(b)作图分析（超大错区）：}题目要求画的是“铅球\textbf{对地面}的压力”，这代表该压力的受力物体是“地面”而不是铅球本身！所以该力的\textbf{作用原点必须死死钉在地面上（也就是球与地面的接触交界切点处）}，不能习惯性画在铅球内部的几何重心上。既然是施加在地面的压迫力，方向自然是垂直于地面长驱直入往下。
3. \textbf{图(c)作图分析（陷阱错区）：}题目要求画的是“电灯\textbf{对电线}的往下拽拉力”，此力的受力体是“电线”！因此，它的发力发源作用点必须只点在悬挂那根笔直电线上（确切而言，即电线的最底末端、与电灯的接触悬挂扣处），不能画在下方电灯本体的内部重心上。同时这股外力是电灯本体对电线施加造成的下方重型坠拉效应，故其方向严格笔直向下。
}

\vspace{0.6cm}

{\small\color{blue!70!black}
\textbf{绘图基本原理与步骤技巧：}\\
\textbf{1. 剥离内秉重心，精确制导局部力的空间定位锚点：}这三张同系列关联图形在应试制卷中的最大反向考量要求就是严格禁止再套用“重心万能论”去强行霸占统筹作图体系。这也是对机器底层坐标系提出了彻底跳关脱骨层级的高频排版验证挑战：我在作图框架里完全禁绝了所有默认的球套圆心算法模块。在 (b) 图中，直接向强制渲染台通过立项设立定触切点 \texttt{(ContactP) at (0, 0)} ，把这张图仅有的外在发力源的坐标点一脚油门死死锁死敲定在了地板接触线上；在 (c) 图中，更是通过重新定义挂载点 \texttt{(WireEnd) at (0, 0.8)} 而直接严密防守截击住了绳系交界点的物理隔离界缝段。通过这种冷酷强硬地解耦所有被分析对象的繁重几何形状与虚无体量、果决改采纯净二维横纵坐标硬编码独立寻址空间节点的形式；把一切物理作用力的起发端精准、粗暴严密地全盘驱赶出了物体的原本内壁正腔，真正向着挑剔严苛的特级中考评卷教师投身输送了出当今数字绘本领域最具霸权统治感、零误伤判正率标兵的一组防杠精严密教学模板图！\\
\textbf{2. 行使强制防御型矢量下斩遮挡图层穿透显印模型：}即使在诸如 (c) 图此类附含多种繁中套繁交叠内切曲弧结构（如大跨口倒锥碗伞罩外嵌深插小弧度灯泡玻璃腹胆轮廓）的经典“旧式欧式吊线垂挂灯”混合拼装包覆组件体壳设计中；我的制图骨架依然清爽游刃有余。由于该反击向下拉扯的绳端作用点恰巧十分凶险刁钻地深刺紧扣在整个灯件外壳的最巅峰防线中心处帽沿中轴端头部位处，且这根粗壮的代表力量的向量拉直箭矢又偏偏被勒令强硬破甲直接冲下笔直猛降扎入其中；为了防止线条被同色填充淹没甚至穿模污染吞噬失位成残图断箭……这根带有剧烈攻击性与贯透刺刺属相的 \texttt{draw[->, ...]} 巨体墨色超黑加粗强拉力线段段块，被我在绘码内部极具针对性设计成“保留后置出场待机排序阵列”——即使在灯具的顶层防漏三角罩跟底封水波灯胆全部被强涂了 \texttt{fill=white} 这霸道具有厚漆白粉无情覆膜隔离特效遮掩渲染罩的前提下：也只凭放任依赖最终行的那条重力出击黑锋针段，就可以利用了在整个庞杂绘板全部渲染指令宣告执行死库耗尽结束之前的最后一截断秒收尾倒数时间流转刹拉中猛烈击穿打印呈现！正是利用直接借调了基于最纯正的 LaTeX TikZ \TeX\ 图层宏大工业引擎自带底蕴时序叠印堆载输出逻辑特质法则所赋予的、雷打不断的神圣“上层栈道盖印高覆下”特权级时空显示遮挡铁律：这根强力下沉式硬刺才得以在这张密闭繁杂架构草图体系里，极不受干扰、冷测明净耀目地当场力战强杀所有内部干扰涂色带防线实现穿刺印透效果；直接直接的无理地达成了手工黑板画老师需要苦心孤诣费劲用粉笔涂抹勾抹才能偶尔避让勉强保全的终极悬线反拽向隐性强扯坠拉物理场景意象，全时标明了代码底层画盘工具相较一切徒手外力所具备拥有着极无边际上限潜力的神迹简化操作碾压力！
}

%% ==============================
%% 第10题与第11题（新）：不同运动状态与附着面下物体重力方向的探究
%% ==============================


\section{解答：不同运动状态与附着面下重力方向探究}

\vspace{0.6cm}

\begin{center}
\begin{tikzpicture}[>=Stealth, line join=round, line cap=round, scale=1.2]

  % ---------------------------------------------------
  % (第10题)：摆动过程中的小球受重力
  % ---------------------------------------------------
  \begin{scope}[shift={(-3.0, 0)}]
    % 天花板
    \draw[black, line width=1.0pt] (-0.8, 2.5) -- (0.8, 2.5);
    \foreach \x in {-0.6, -0.4, ..., 0.6} {
      \draw[black, line width=0.6pt] (\x, 2.5) -- (\x+0.12, 2.65);
    }
    
    % 悬挂点 O 和 纯细刚绳子
    \coordinate (O) at (0, 2.5);
    \coordinate (BallCenter) at (0.8, 0.6);
    % 为了图层防重叠，直接让随后渲染的球体覆盖掉绳子下端交接突刺部位
    \draw[black, line width=1.2pt] (O) -- (BallCenter); 
    
    % 主体小球 (覆白底掩饰绳段穿模)
    \draw[black, line width=1.2pt, fill=white] (BallCenter) circle (0.18);
    
    % 运动方向虚线动态箭头尾流 (用贝塞尔曲线平滑绘制向左摆动的轨迹)
    \draw[<-, dashed, gray!80!black, line width=1.2pt] (0.0, 0.3) .. controls (0.3, 0.2) and (0.6, 0.3) .. (1.0, 0.5);
    
    % 核心踩分点：重力作图 (不管怎么摆，必须竖直向下)
    % 借由题干常数直接算出重力：G = mg = 2kg * 10N/kg = 20N
    \fill[black] (BallCenter) circle (1.5pt);
    \draw[->, black, line width=1.8pt] (BallCenter) -- ++(0, -1.6) node[right, font=\large\bfseries] {$G = 20\,\text{N}$};

    \node at (0, -1.8) {\Large (第$10$题)};
  \end{scope}

  % ---------------------------------------------------
  % (第11题)：水平平原面A与陡峭斜面B的重力对比
  % ---------------------------------------------------
  \begin{scope}[shift={(3.0, 0)}]
    % 综合相连大地（左侧平铺水平，右侧拔起斜坡）
    % 平铺水平底面
    \draw[black, line width=1.0pt] (-2.5, 0) -- (0, 0);
    \foreach \x in {-2.4, -2.2, ..., -0.2} {
      \draw[black, line width=0.6pt] (\x, 0) -- (\x-0.12, -0.12);
    }
    
    % 45° 陡峭斜坡面
    \draw[black, line width=1.0pt] (0, 0) -- (1.8, 1.8);
    \foreach \x/\y in {0.1/0.1, 0.3/0.3, 0.5/0.5, 0.7/0.7, 0.9/0.9, 1.1/1.1, 1.3/1.3, 1.5/1.5, 1.7/1.7} {
      % 法线正交角度强行切除哈希底纹纹理标刺 (+0.12, -0.12)
      \draw[black, line width=0.6pt] (\x, \y) -- (\x+0.12, \y-0.12);
    }
    
    % 物体 A (平躺在水平面上)
    \draw[black, line width=1.2pt, fill=white] (-1.6, 0) rectangle (-0.6, 0.6);
    \node[font=\large, yshift=0.5pt] at (-1.1, 0.3) {A};
    \coordinate (CenterA) at (-1.1, 0.3);
    
    % A 本分的重力 (直接从四方 A 箱心朝下垂落)
    \fill[black] (CenterA) circle (1.5pt);
    \draw[->, black, line width=1.8pt] (CenterA) -- ++(0, -1.2) node[right, font=\large\bfseries] {$G$};
    
    % 物体 B (依赖在斜面上)
    % 调用 rotate=45 局部环境，使得画法大为简化
    \begin{scope}[shift={(0.8, 0.8)}, rotate=45]
      \draw[black, line width=1.2pt, fill=white] (-0.5, 0) rectangle (0.5, 0.6);
      % 强行矫正视线！让内部字母 B 保持原本竖直状态，不随方块自体倾倒
      \node[font=\large, rotate=-45, yshift=-0.5pt] at (0, 0.3) {B};
      
      % 【关键信标机制】
      % 在这个会诱导向量倾斜的 45度密闭小黑屋内，只保存定焦坐标信标而不绘图！
      \coordinate (CenterB_local) at (0, 0.3);
    \end{scope}
    
    % *生死攸关的跳出域渲染*
    % B 的重力必须生生退逃简化跳离出上一个 [rotate=45] 的区域法阵，从根本上逃离倾斜面力场污染。
    % 回归世界真源坐标系后，再取出 CenterB_local 纵向拉出竖直重力直线箭矢！
    \fill[black] (CenterB_local) circle (1.5pt);
    \draw[->, black, line width=1.8pt] (CenterB_local) -- ++(0, -1.2) node[right, font=\large\bfseries] {$G$};

    \node at (0, -1.8) {\Large (第$11$题)};
  \end{scope}

\end{tikzpicture}
\end{center}

\vspace{0.6cm}

{\small\color{gray}
\textbf{解题说明：}\\
这组极具防偏对比意义的试题唯一的终极考点建立在一个物理公理上：\textbf{重力的指向，永远是保持严厉竖直向下的（垂直于水平面向下方）}。这一铁定则不受制于物体本身正在经历的任何复杂运动（例如向左剧烈拉扯摆动的第 10 题），也不受制于它所附着支撑大平面的接触面斜率扭曲度（例如处于高斜大陡坡上的第 11 题物块 B），不可被视觉假象产生诱拐偏滑。\\
1. \textbf{第 10 题作图破障分析：}许多考生因为视觉残留的向左大摆动虚线和有倾斜角度的硬拉拉绳的强烈诱捕，导致顺手将重力也给画向侧方。真正的满分逻辑是：当图面所有环境都在宣称处于激烈变化时，握住那定海神针；将起画点钉死于小球重心，直截了当地拉一条向下的纯重锤线。同时利用体贴显露的常量参配（$m=2\,\text{kg}$）迅速写明重力推算值：$G = mg = 2\,\text{kg} \times 10\,\text{N/kg} = 20\,\text{N}$。\\
2. \textbf{第 11 题作图拆阵分析：}它旨在用右面在夸张倾倒山坡上的物块 B 去诱使初学者们翻车——大批考生会习惯性顺势把 B 表面往下延伸出的重力大线给画偏，画成了平行亦或是垂直于高斜面这致命大忌。此处的解题戒律是：全视界只唯一下方大白底最准纸边作为物理校验对照长线，强无视掉所有的陡坡框选干扰框架，向底栏直下不妥协的拉一记最长向下的强垂斩落。
}

\vspace{0.6cm}

{\small\color{blue!70!black}
\textbf{制图（代码）防倾覆防穿刺高阶处理解密：}\\
\textbf{1. 非规动态伪拉平顺虚抛流控（贝塞尔三阶）：}为了高灵敏度指示还原那根用于表征此际此刻主动态正在顺延向着侧方狂奔激荡的大虚线后拽跟踪轨迹；我摒弃了完全呆板生硬的预设图例系统死画库或者费心耗力多角拼接算法；当场切调入那富含极致数理顺畅魅力的贝塞尔三阶曲线控制拉引流：\texttt{.. controls () and () ..} ；仅仅给定初末二落点搭配悬空预控引导的双锚标定位参数化诱发，编译即会自动算抛渲染出极致饱满吻切主球行径弧线的唯美运迹线条动态走势。\\
\textbf{2. （核爆级制图大招）降域内信标封存，跳出层绝离渲染斩击抗旋保护体防御法：}针对绘制悬爬在最受歪向斜倒病毒感染重灾区（斜角高倾位 $45^\circ$ 全重转法境）那侧倒斜趴大箱子 B 时，我迫不得已动用了\texttt{\textbackslash begin\{scope\}[rotate=45]} 来进行省脑算量的作画框倒模复写生成；但假如在里面就执行那神圣重力垂落动作 \texttt{-{}- ++(0,-1.2)} 的话必定会在计算时遭系统 $45^\circ$ 图层斜向同调吞噬同化扭偏画成斜伪重力线。为保其金刚纯血原向轴正交不歪斜，我极限强令所有关于指向的操作指令全部逃离闭环封死空间，仅在此法域地底内烙写一根单纯引爆定子座印大坐标：\texttt{\textbackslash coordinate (CenterB\_local)...} ，随即果绝一刀 \texttt{\textbackslash end\{scope\}} 切断跳回大界本源平原场后；再次提出那一安插的定点信物引雷爆点向着大一统本源纯向施加强令输出力：直接画出绝笔长垂死硬的大地引力纯垂直黑箭！这可以说是一招应对后续全部局部偏离干扰法框硬核无敌的神图防御拔剑制底招。
}

%% ==============================
%% 第13题：探究弹簧测力计原理及刻度板绘制
%% ==============================


\section{解答：探究弹簧测力计原理及刻度板绘制}

\vspace{0.6cm}

\textbf{（1）从表中的数据可以看出弹簧的伸长量 $\Delta l$ 与拉力 $F$ 的关系为：}\\
\underline{\quad 在弹性限度内，弹簧的伸长量与受到的拉力成正比。\quad}

\vspace{0.4cm}

\textbf{（2）小明继续实验，得到如下数据：} \\
从这次数据看出：拉力达到 \underline{\quad $\textbf{11}$ \quad} N 时，弹簧的伸长量和拉力的关系就改变了。因此，该弹簧测力计的最大测量值为 \underline{\quad $\textbf{10}$ \quad} N。

\vspace{0.4cm}

\textbf{（3）请你根据表中的数据，在图中用刻度尺画出这个弹簧测力计的刻度板（分度值为 1N）。}

\vspace{0.4cm}

\begin{center}
\begin{tikzpicture}[>=Stealth, line join=round, line cap=round, scale=1.0]
  % ---------------------------------------------------
  % 1. 测力计面板外轮廓 (含上下半圆凸起，且加宽边距防数字重叠)
  % ---------------------------------------------------
  \draw[black, line width=1.0pt, fill=white] 
    (-1.0, 0.8) -- (-0.35, 0.8) arc(180:0:0.35) -- (1.0, 0.8) -- 
    (1.0, -4.6) -- (0.35, -4.6) arc(360:180:0.35) -- (-1.0, -4.6) -- cycle;
    
  % ---------------------------------------------------
  % 2. 中间的直槽 (上下两端圆润过渡)
  % ---------------------------------------------------
  \draw[black, line width=0.8pt, fill=gray!5] 
    (-0.15, -4.3) -- (-0.15, 0.3) arc(180:0:0.15) -- 
    (0.15, -4.3) arc(360:180:0.15) -- cycle;
    
  % ---------------------------------------------------
  % 3. 顶部的单位 'N' 
  % ---------------------------------------------------
  \node[font=\large] at (0, 0.9) {N}; 
  
  % ---------------------------------------------------
  % 4. 绘制 0 刻度线 (直接从右侧槽边伸出，高度契合原图外壳)
  % ---------------------------------------------------
  \draw[black, line width=1.0pt] (0.15, 0) -- (0.4, 0); 
  \node[right, font=\small, inner sep=2pt] at (0.4, 0) {0};

  % ---------------------------------------------------
  % 5. 严格按物理极距 (0.4cm) 补充的蓝色答题层刻度
  % ---------------------------------------------------
  \foreach \F in {1, ..., 10} {
    \pgfmathsetmacro{\y}{-0.4 * \F}
    % 画分度值为 1N 的主刻度线
    \draw[blue!80!black, line width=1.0pt] (0.15, \y) -- (0.4, \y);
    \node[right, text=blue!80!black, font=\small, inner sep=2pt] at (0.4, \y) {\F};
  }
\end{tikzpicture}
\end{center}

\vspace{0.6cm}
\hrule
\vspace{0.4cm}

{\small\color{gray}
\textbf{解题分析与说明：}\\
1. \textbf{定性归纳第一题：}在 $1\sim 7\,\text{N}$ 的区间内，拉力 $F$ 每增加 $1\,\text{N}$，伸长量 $\Delta l$ 恒定增加 $0.40\,\text{cm}$。这验证了胡克定律的经典表述：“在弹性限度内，弹簧的伸长量与受到的拉力成正比”。\\
2. \textbf{寻找拐点破第二题：}拉力到 $8, 9, 10\,\text{N}$ 时，其伸长量 $3.20, 3.60, 4.00\,\text{cm}$ 依然吻合正比例。但当拉力达到 \textbf{$11\,\text{N}$} 时，伸长量实测值变为 $4.50\,\text{cm}$（理论应为 $4.40\,\text{cm}$），说明正比例关系被彻底打破。因此，测量上限必须卡死在正比区间的最高边缘也就是 \textbf{$10\,\text{N}$}。\\
3. \textbf{严谨绘制标线：}明确了最大量程为 $10\,\text{N}$，且分度值为 $1\,\text{N}$，就意味着需要等比例精确无误地画出 $10$ 根刻度线。每根线的间隔长度必须严格依据实测物理极距 $0.40\,\text{cm}$。
}

\vspace{0.6cm}

{\small\color{blue!70!black}
\textbf{代码自动化处理解析：}\\
\textbf{1. 物理厘米级等距映射刻度：}不同于手绘粗略示意，代码可以直接绑定物理真实步距。在默认 $1.0 = 1.0\,\text{cm}$ 的比例系中，循环逻辑 $\texttt{\textbackslash pgfmathsetmacro\{\textbackslash y\}\{-0.4 * \textbackslash F\}}$ 驱动画笔，确保在纸面上呈现出的真实物理 $0.4\,\text{cm}$ 间隔极距跨步，还原了实验中游标下沉的物理真实距离。\\
\textbf{2. 循环量程控制锁：}刻度条数的决定并非随意排布，而是由题干第二问的“弹性限度安全区域”强力关联绑定。\texttt{\textbackslash foreach} 循环被强制封死在最大值 \texttt{10} 处，阻断了刻度越界超画的安全隐患。\\
\textbf{3. 答题区色彩直观防线：}为了突显被学生后期补绘卷面充填出来的内容，我专门动用了 \texttt{blue!80!black} 为所有新生刻度线和标识数字进行深度强着色。这不仅展现了答题前后的步骤层叠感，更使正确刻划痕变得格外醒目。
}

%% ==============================
%% 附加图：m-V (质量-体积) 图像基础网格
%% ==============================

